% Appendix D: the catalogue (how it was built and checked; the complete tables).
\section{The catalogue}\label{app:catalogue}

The atlas of Exposé~\ref{sec:atlas} rests on a catalogue of \nMethods{} methods and \nPapers{} papers. This appendix
says how the catalogue was assembled and fact-checked, how each method received its seven coordinates, its arrows and
its obstructions, and how these were checked in turn. It ends with the complete tables: every method with its
coordinates and its reference (Table~\ref{tab:catalogue}), every recorded arrow with its map $h$
(Table~\ref{tab:arrows}), and the arrows and obstructions counted by rule and by test
(Table~\ref{tab:app-rules-tests}).

\paragraph{Data files} The data accompany this paper in machine-readable form, generated from the same source as the tables
below and as the companion site.
\begin{description}[leftmargin=1.2em,itemsep=2pt]
  \item[\texttt{catalogue.json}] Every method (\nMethods{}) and every paper (\nPapers{}), each with its verified
  metadata (title, full author list, year of the first arXiv version, venue, links to the paper and the code) and its
  BibTeX record. A method entry adds the formula, the trainable-parameter count, the locus, the modification kind, the
  merge verdict with its reason, the initialisation, the relations to other methods, the sources consulted, the seven
  coordinates with a justification for each, the categorical reading, and the method's arrows and obstructions. The
  file also holds the definitions of the axis values, with the question that decides each axis, and the arrow rules.
  \item[\texttt{catalogue.csv}] One row per method, with the main fields of \texttt{catalogue.json} and one column per
  coordinate.
  \item[\texttt{arrows.csv}] Every recorded arrow $M\to N$ with its rule, its map $h$ and a note that names the instance
  of $N$ it lands in (rank, pointing or regime).
  \item[\texttt{obstructions.csv}] Every recorded obstruction $M\not\to N$ with the test that rules the arrow out and
  the invariant that fails.
  \item[\texttt{propositions.json}] Every numbered result (\nPropositions{} of them) with its statement, proof, review
  status, second-round verdict, public statement and referee notes.
  \item[\texttt{framework\_checks.py}] The numerical checks of Appendix~\ref{app:repro}.
\end{description}

\subsection{Families and sources}\label{app:catalogue-families}

Every entry was written once and then attacked once, by a separate reader looking for what the first got wrong. The
literature was split into the thirteen families of Table~\ref{tab:families}, then completed by four gap-filling
batches, which added the methods a completeness review found missing, and by two addenda on 2 October 2026: methods the authors reported missing, and a sweep of the 2026 literature run in two passes
(\S\ref{app:catalogue-2026}). The family boundaries owe much to earlier maps of the field: the unified view of
\citet{he2021unified}, the delta-tuning study of \citet{ding2022parameter}, the guide of \citet{lialin2023scaling}, the
survey of \citet{han2024parameter} and the modular view of \citet{pfeiffer2023modular}.

\begin{table}[t]
\centering\footnotesize
\caption{How the catalogue was assembled. Each family was written by one reader and attacked by another; the last column
(Corr.) counts the logged corrections. The families overlap, so the records were merged into the catalogue on a canonical
identifier.}\label{tab:families}
\begin{tabularx}{\linewidth}{@{}>{\raggedright\arraybackslash}p{4.1cm}>{\raggedright\arraybackslash}Xrrr@{}}
\toprule\headrow
\textbf{Family} & \textbf{What it covers} & \textbf{Methods} & \textbf{Papers} & \textbf{Corr.}\\\midrule
LoRA core & LoRA and its variants & 20 & 12 & 21\\
Initialisation and rank & Base-point splits; rank allocation & 25 & 6 & 46\\
Shape & Kronecker, Hadamard, tensor and sparse shapes & 29 & 9 & 49\\
Multiplicative & Group and monoid actions on $W_0$ & 32 & 9 & 54\\
Sharing & Shared or random bases & 28 & 4 & 46\\
Selective & Subsets of existing weights & 26 & 7 & 37\\
Adapters & Inserted modules and side networks & 31 & 10 & 30\\
Prompts & Soft prompts, prefixes and learned states & 27 & 13 & 39\\
Quantisation and memory & Quantised bases, gradient memory & 24 & 5 & 44\\
Composition & Merging, routing, mixtures & 42 & 8 & 61\\
HF PEFT & The library's tuners and utilities & 74 & 10 & 53\\
Frontier & 2025--2026: RL, reasoning, diffusion & 21 & 5 & 48\\
Theory & Theory and foundations & 0 & 50 & 23\\
Gap-filling (4 batches) & Methods a completeness review found missing & 36 & 18 & 62\\
Addendum (2 Oct 2026) & Methods the authors reported missing & 1 & 0 & 2\\
Addendum (2026 sweep, pass 1) & arXiv, January to September 2026; ICLR and ICML 2026 & 29 & 8 & 126\\
Addendum (2026 sweep, pass 2) & Further candidates, mostly ICLR, ICML and NeurIPS 2026 & 32 & 6 & 164\\\midrule
\textbf{All records} & Before merging & \textbf{477} & \textbf{180} & \textbf{905}\\
\textbf{Catalogue} & After merging duplicates & \textbf{\nMethods} & \textbf{\nPapers} & \\
\bottomrule
\end{tabularx}
\end{table}

\paragraph{Merging}
The families overlap, so the $477$ method records and $180$ paper records were merged into \nMethods{} methods and
\nPapers{} papers. Records were matched on a canonical identifier and a table of aliases, because matching by name would
fold LoRA+ \citep{hayou2024lora} into LoRA \citep{hu2021lora} and O-LoRA \citep{wang2023orthogonal} into OLoRA
\citep{buyukakyuz2024olora}. Of the \nMethods{} entries, $358$ are methods proper; the others are $17$ training recipes,
$4$ software systems, $2$ analyses and $1$ baseline (linear probing).
The first arXiv versions of the methods run from 2013 to 2026, and $47$ methods were first posted in 2026.

\paragraph{The bibliography}
The BibTeX record of each entry comes from the arXiv metadata of its identifier, with the full author list, for $493$
identifiers, and was written by hand for the nine entries without one. These $502$ records
form the catalogue part of this paper's bibliography, and the citation keys of Table~\ref{tab:catalogue} point into it.

\subsection{Fact-checking}\label{app:catalogue-verification}

Each field was checked against a primary source.
\begin{itemize}
  \item Titles, authors and first-version years were checked against the arXiv metadata.
  \item Venues were checked against the official records: OpenReview, the ACL Anthology, PMLR and CVF open access, with
  Crossref, Semantic Scholar and DBLP for the remaining cases.
  \item Formulas and parameter counts were checked against the papers, and against the official code where an
  initialisation was in doubt.
  \item Every claim about Hugging Face PEFT \citep{xmangrulkar2022peft} was checked against its source at commit
  \texttt{85be3f3} (\texttt{85be3f3a1a58c4d153ffc4631c57196430af87a4}) of 30 September 2026. The library's code
  settled many questions of convention, such as the side of $W$ on which a rotation acts, the default initialisation of
  each factor, and whether a scale such as $\alpha/r$ is applied. The column ``HF'' of Table~\ref{tab:catalogue} marks
  the \nHF{} methods available at that commit.
\end{itemize}
Every change was logged with its reason, $905$ in all (Table~\ref{tab:families}). The merge verdict of each method is
recorded beside its reason: $250$ methods fold back into the frozen weights always, $41$ under stated conditions and $91$
not at all. This verdict concerns the deployed method; the merge coordinate of Exposé~\ref{sec:atlas} refines it by the
rewrite rules of \thmref{Theorem}{VI.3}.

\subsection{The 2026 sweep and the authors' addendum}\label{app:catalogue-2026}

On 2 October 2026 the authors pointed out methods missing from the catalogue. LOFT \citep{zhao2026loft} was added at once as an addendum of its own,
checked against arXiv and the paper (two corrections logged), and classified like every other method.
Its coordinates are R, Orb, regular, none, $\CO\times\CO$, M1 and idempotent; it has arrows into OFT, LoRA-XS, PSOFT
and LoRA, and for $r\ge4$ its first-order image rules out a smooth pointed arrow into HRA at HF's paired
initialisation (test O3). Two
independent referees checked these coordinates, arrows and the obstruction, numerically and against the paper and the
code. \S\ref{sec:erlangen-loft} reads LOFT as a fixed leaf of HRA's meet.

The other gaps were sought by a sweep of the 2026 literature, run in two passes on the same day. A pool of $265$ candidates was
drawn from six slices:
\begin{itemize}
  \item three quarterly arXiv slices, for first versions posted from January to March, April to June and July to
  September 2026, with $60$ candidates each;
  \item the accepted papers of ICLR 2026, $5{,}468$ of them screened through their abstracts, with $48$ candidates;
  \item the accepted papers of ICML 2026, $6{,}797$ screened, with a second pass over ACL, EACL and CVPR 2026, together
  giving $31$ candidates;
  \item the 2026 papers citing anchor methods (LoRA, DoRA, PiSSA, QLoRA, OFT, VeRA, HRA, GaLore, PSOFT and LoRA-GA),
  ranked by citations, recency and acceptance, with the 2026 release notes and pull requests of Hugging Face PEFT,
  giving $6$ candidates.
\end{itemize}
Acceptance at ICLR and ICML 2026 was confirmed by an exact title match against the official lists of accepted papers
or by the authors' own arXiv comment. Candidates were matched against the catalogue by arXiv identifier and by name, so that near-homonyms stay
apart: the entropy-guided FlexLoRA of 2026 \citep{liu2026flexlora} is not the federated FlexLoRA of 2024
(arXiv:2402.11505), and the per-token MoLoRA of 2026 (arXiv:2603.15965) is not the MoLoRA of \citet{zadouri2023pushing}.
The first pass admitted $29$ methods and $8$ papers. The second re-screened the rest of the pool and admitted $32$ methods
and $6$ papers, mostly from ICLR, ICML and NeurIPS 2026. Each admitted entry went through the same fact-checking as the
families, with $126$ and $164$ logged corrections, and each of the $61$ new methods received an independent check of its
coordinates and arrows on the same day. The arXiv queries searched abstracts for LoRA, PEFT, adapter,
parameter-efficient and prompt tuning, so a 2026 method paper whose abstract uses none of these terms lies outside the
sweep's reach.

\subsection{Coordinates, arrows and obstructions}\label{app:catalogue-classification}

\paragraph{Assignment}
Each method was classified from its formula, in thirteen batches: ten of $32$ methods for the catalogue as it stood
on 1 October 2026, one for the authors' addendum, and one for each pass of the 2026 sweep ($29$ and $32$ methods). A first reader answered the
seven questions of \S\ref{sec:atlas-axes}, which fix the kind, the shape of the image germ, the base point, the gauge,
the covariance group, the merge type and the rebasing closure, following the decision procedures of
\S\ref{sec:atlas-decide} and the conventions of \S\ref{sec:atlas-conventions}. Each answer carries a short
justification, kept in \texttt{catalogue.json}. A null coordinate, printed ``--'', means that the axis does not apply,
for instance the shape of a method with no germ at $\theta_0$. Where an answer turns on the dimension of an image or a
fibre, or on the first-order image at the initialisation, it was confirmed by a float64 Jacobian rank at a random point
and at the base point (\bgref{C.3}).

Each arrow $M\to N$ comes with an explicit pointed smooth map $h$ satisfying $\rho_N\circ h=\rho_M$, obtained by one of
the rules A2--A5 of \S\ref{sec:atlas-rules}, and with the instance of $N$ it lands in: a rank, a pointing or a
hyperparameter regime. Each obstruction $M\not\to N$ names the test O1--O3 of \S\ref{sec:atlas-tests} that rules the
arrow out and the invariant that fails. The interval arrows $\text{Frozen}\to M\to\text{Full FT}$ of rule A1 hold for
every weight-space method and are not recorded. The \nCore{} core methods are those named in the periodic table
(Table~\ref{tab:periodic}) together with a curated set of others; the flag is not a coordinate.

\paragraph{Independent check}
A second reader re-derived every coordinate from the formula and re-checked every arrow, by recomputing
$\rho_N\circ h$ and the pointing $h(q_0^M)$, and every obstruction, by recomputing the failing invariant. Each method
carries a check note with the outcome. More than half of the notes record no change; the others record each change with
its reason, and $47$ record a numerical Jacobian computation (a rank, a fibre dimension or a deficiency).

A consolidation pass on 1 October 2026, over the $320$ methods classified by then, unified the data and touched $135$
methods, each change logged in its check note.
\begin{enumerate}[label=(\arabic*)]
  \item $137$ coordinates used a value ``n/a'' that the axis definitions do not contain. $136$ became null, matching
  the ``--'' of the tables, and the base point of GenKnowSub became ``unpointed'', as for the other parameterless
  post-hoc operations. The merge axis keeps its legacy key \texttt{na}, read as null.
  \item The row labels of the periodic table were unified across batches.
  \item The covariance of LISA and BAdam changed from $\Pi$ to $\GL\times\GL$, matching HiFT, surgical fine-tuning and
  BitFit: each stage frees whole boxes, and the layer partition breaks no neuron permutation.
  \item Four obstruction reasons now say which instance of the target they concern (DSEE, GaLore, Spectral Adapter$^A$
  and MoV, each against LoRA).
\end{enumerate}
Where a formula disagreed with an earlier draft of the periodic table, the value the formula supports was kept and the
table of Exposé~\ref{sec:atlas} was revised to agree, in thirteen places. For example, LST, UniPELT
and MAM are unpointed extensions, MPT and Compacter have torus gauges, the Houlsby adapter's GELU or swish leaves only the
permutations of hidden units, and the gauges of DeLoRA and AdaLoRA are non-linear. The $62$ later entries, LOFT and the
$61$ methods of the 2026 sweep, were classified and checked in the same way.

\paragraph{Structural checks on the final data}
For this appendix the structural checks were re-run on the ancillary files. Every arrow and every obstruction joins two
catalogued methods and names a valid rule or test; every arrow has a map $h$ and every obstruction a reason; every
coordinate is a value defined for its axis; and no arrow or obstruction is recorded twice. A check at the level of method
identifiers flags four kinds of pattern, and each was resolved by the instance it concerns.
\begin{itemize}
  \item \emph{Pairs with both an arrow and an obstruction} ($26$). In every case the two concern different instances of
  the target: a different rank, pointing, regime or variant (\S\ref{sec:atlas-checks}). For example, ABBA is isomorphic
  to LoHa pointed at the factors of $[W_0]_{r_1}$, while the O3 obstruction concerns LoHa at HF's random default
  pointing.
  \item \emph{Arrows from a method to itself} ($8$). Each is an inclusion between instances of one family, such as
  LoRA$_r\hookrightarrow{}$LoRA$_{r+1}$ or ReLoRA$^K\hookrightarrow{}$ReLoRA$^{K+1}$.
  \item \emph{Cycles} ($13$ groups of mutually reachable methods). The cycles in these groups are of four kinds.
  Eight are isomorphisms: DoRA and BiDoRA, ReLoRA and COLA, AsyLoRA and LoRA-FA, AdaptFormer and the parallel adapter,
  S-LoRA and Punica, EFT and proxy tuning, LoRAMoE and MoLoRA, and AdaMerging and LoRA soups (CAT), which differ only in
  pointing. Three join methods with equal images that are not isomorphic: prompt tuning with MPT and with P-tuning, and
  QLoRA with IR-QLoRA. Five go up in rank or in the number of stages and are not isomorphisms: LoRA with ReLoRA, LoRA
  with PEANuT, LoRA-FA with Flora, surgical fine-tuning with HiFT, and KronA with LoKr. Initial-state tuning and prefix
  tuning simulate each other on S4 state-space layers only.
  \item \emph{Obstructions $M\not\to N$ with no direct arrow but a path of arrows from $M$ to $N$} ($20$). In each,
  the obstruction concerns a fixed rank or
  pointing of the target and the path ends in another instance. FourierFT with $n_c$ frequencies, for example, maps
  into LoRA-FA$_{2n_c}$ and so into LoRA$_{2n_c}$, while the O2 obstruction concerns LoRA$_r$ at a matched budget.
\end{itemize}
No arrow contradicts an obstruction once the stated rank or pointing is taken into account. The eight places where the
data depart from the conventions of \S\ref{sec:atlas-conventions} are listed in \S\ref{sec:atlas-checks}, and
Table~\ref{tab:atlas-counts} counts the methods per coordinate value.

\begin{table}[t]
\centering\footnotesize
\caption{The \nArrows{} recorded arrows by rule (top) and the \nObstructions{} recorded obstructions by test (bottom).
The rules and tests are those of \S\ref{sec:atlas-arrows}; $161$ of the arrows join two core methods. Every arrow is
listed in Table~\ref{tab:arrows}, and every obstruction, with its failing invariant, in the ancillary file
\texttt{obstructions.csv}.}\label{tab:app-rules-tests}
\begin{tabularx}{\linewidth}{@{}lrX@{}}\toprule\headrow
\textbf{Rule} & \textbf{Count} & \textbf{How $h$ arises}\\\midrule
A2 & 61 & freezing a block of coordinates at its base value, $h(q')=(q',q_0'')$\\
A3 & 329 & an explicit smooth pointed factorisation $\rho_M=\rho_N\circ h$\\
A4 & 2 & a projection out of a fibre product\\
A5 & 11 & a sum or an added stage, $h(q)=(q,q_0^N)$\\
\bottomrule\end{tabularx}

\medskip
% generated by tools/make_tables.py
\begin{tabular}{@{}lrl@{}}\toprule\headrow\textbf{Test} & \textbf{Count} & \textbf{What fails}\\\midrule
O1 & 288 & an image invariant (Gram matrix, spectrum, row lines, zero pattern) or image inclusion\\
O2 & 129 & dimension or cut-rank\\
O3 & 68 & first-order image or base-point stratum\\
\bottomrule\end{tabular}

\end{table}

\subsection{The complete tables}\label{app:catalogue-tables}

Table~\ref{tab:catalogue} lists every method with its reference, the year of its first arXiv version and its seven
coordinates. A method without an arXiv version carries the year it first became public. The year in a citation is that
of the version cited, so the two can differ: LoRA-S, for instance, first appeared in 2025 and is cited as
\citet{zheng2025lora}, its ICLR 2026 version. The rows are grouped by kind, with the reparametrisations split as in the periodic table (additive, action,
Hadamard and selective), and sorted by name within each group. The abbreviations are those of
Table~\ref{tab:atlas-counts}, with ``unpt.''\ for unpointed, ``nonlin.''\ for a non-linear gauge, GL$\times$T for
$\GL$ times a Hadamard torus, O$\times$O for the spectral-frame value $\CO\times\CO$, ``idemp.''\ for idempotent and
``sched.''\ for a scheme that is itself a rebasing schedule. The merge value \texttt{na} is the legacy key for
backward methods, read as ``--''.

Table~\ref{tab:arrows} lists every recorded arrow, sorted by source, with its rule and its map $h$. Each $h$ is stated
in the notation of the target's formula in \texttt{catalogue.json}, and the instance of the target (its rank, its
pointing or its regime) is in the note column of \texttt{arrows.csv}.

% generated by tools/make_tables.py
\begingroup\scriptsize\setlength{\tabcolsep}{2.1pt}\renewcommand{\arraystretch}{1.08}
\begin{longtable}{@{}>{\raggedright\arraybackslash}p{2.53cm}>{\raggedright\arraybackslash}p{3.2cm}r l l l l l l l l@{}}
\caption{The complete catalogue: every method with its seven coordinates (\S\ref{sec:atlas}). Core methods are marked $\bullet$; HF marks methods available in Hugging Face PEFT at the pinned commit. ``--'' means the axis does not apply.}\label{tab:catalogue}\\
\toprule\headrow\textbf{Method} & \textbf{Reference} & \textbf{Year} & \textbf{Kind} & \textbf{Shape} & \textbf{Base pt.} & \textbf{Gauge} & \textbf{Covar.} & \textbf{Merge} & \textbf{Rebase} & \textbf{HF}\\\midrule
\endfirsthead
\multicolumn{11}{@{}l}{\footnotesize\itshape Table~\ref{tab:catalogue}, continued}\\
\toprule\headrow\textbf{Method} & \textbf{Reference} & \textbf{Year} & \textbf{Kind} & \textbf{Shape} & \textbf{Base pt.} & \textbf{Gauge} & \textbf{Covar.} & \textbf{Merge} & \textbf{Rebase} & \textbf{HF}\\\midrule
\endhead
\midrule
\multicolumn{11}{@{}r}{\footnotesize\itshape continued on the next page}\\
\endfoot
\bottomrule\endlastfoot
\multicolumn{11}{@{}l}{\sffamily\bfseries\color{tide}R, additive}\\[1pt]
ABBA & \citealp{singhal2025abba} & 2025 & R & Cone & apex & GL$\times$T & Mon$\times$Mon & M1 & span & \\
$\bullet$\,AdaLoRA & \citealp{zhang2023adalora} & 2023 & R & Cone & apex & nonlin. & GL$\times$GL & M1 & span & \checkmark\\
$\bullet$\,AdaMSS & \citealp{zheng2025adamss} & 2025 & R & Cone & apex & GL & O$\times$O & M1 & span & \checkmark\\
AdaPaD & \citealp{su2026adapad} & 2026 & R & Cone & defect & GL & GL$\times$GL & M1 & span & \\
ALoRA & \citealp{liu2024alora} & 2024 & R & Cone & apex & GL & GL$\times$GL & M1 & span & \\
AltLoRA & \citealp{yu2025altlora} & 2025 & R & Cone & apex & GL & GL$\times$GL & M1 & span & \\
ApiQ & \citealp{liao2024apiq} & 2024 & R & Cone$^{*}$ & defect & GL & Mon$\times$Mon & Mq & span & \\
AR1-ZO & \citealp{chen2026ar1} & 2026 & R & Cone & apex & GL & GL$\times$GL & M1 & span & \\
$\bullet$\,Astra & \citealp{liu2026astra} & 2026 & R & Cone$^{*}$ & regular & GL & O$\times$O & M1 & span & \checkmark\\
AsyLoRA & \citealp{zhu2024asymmetry} & 2024 & R & Lin & regular & none & frame & M1 & idemp. & \\
AutoLoRA & \citealp{zhang2024autolora} & 2024 & R & Cone & apex & GL & GL$\times$GL & M1 & span & \\
$\bullet$\,BD-LoRA & \citealp{wang2025block} & 2025 & R & Cone & apex & GL & $\Pi$ & M1 & span & \checkmark\\
$\bullet$\,C3A & \citealp{chen2024parameter} & 2024 & R & Lin & regular & none & $\Pi$ & M1 & idemp. & \checkmark\\
Concept Sliders & \citealp{gandikota2023concept} & 2023 & R & Cone & apex & GL & GL$\times$GL & M1 & span & \\
$\bullet$\,CorDA & \citealp{yang2024corda} & 2024 & R & Cone$^{*}$ & regular & GL & O$\times$O & M1 & span & \checkmark\\
$\bullet$\,DEFT & \citealp{kumar2025deft} & 2025 & R & Cone & apex & GL & GL$\times$GL & M1 & span & \checkmark\\
$\bullet$\,DeLoRA & \citealp{bini2025delora} & 2025 & R & Cone & apex & nonlin. & GL$\times$GL & M1 & span & \checkmark\\
DoRA (dynamic rank) & \citealp{mao2024dora} & 2024 & R & Cone & apex & nonlin. & GL$\times$GL & M1 & span & \\
DSEE & \citealp{chen2021dsee} & 2021 & R & Cone & apex & GL & Mon$\times$Mon & M1 & span & \\
$\bullet$\,DyLoRA & \citealp{valipour2022dylora} & 2022 & R & Cone & apex & torus & GL$\times$GL & M1 & span & \\
ElaLoRA & \citealp{chang2025elalora} & 2025 & R & Cone & apex & nonlin. & GL$\times$GL & M1 & span & \\
$\bullet$\,EVA & \citealp{paischer2024parameter} & 2024 & R & Cone & apex & GL & GL$\times$GL & M1 & span & \checkmark\\
$\bullet$\,FacT & \citealp{jie2022fact} & 2022 & R & Cone & apex & GL & $\Pi$ & M1 & span & \\
FedRot-LoRA & \citealp{zhang2026fedrot} & 2026 & R & Cone & apex & GL & GL$\times$GL & M1 & span & \\
$\bullet$\,FFA-LoRA & \citealp{sun2024improving} & 2024 & R & Lin & regular & none & frame & M1 & idemp. & \\
Flat-LoRA & \citealp{li2024flat} & 2024 & R & Cone & apex & GL & GL$\times$GL & M1 & span & \\
FlexLoRA & \citealp{liu2026flexlora} & 2026 & R & Cone & apex & nonlin. & GL$\times$GL & M1 & span & \\
FLoRA (Tucker core) & \citealp{si2024maintaining} & 2024 & R & Cone & apex & GL & GL$\times$GL & M1 & span & \\
FLoRG & \citealp{meng2026florg} & 2026 & R & Cone & defect & GL & frame & M1 & span & \\
$\bullet$\,FourierFT & \citealp{gao2024parameter} & 2024 & R & Lin & regular & none & $\Pi$ & M1 & idemp. & \checkmark\\
FourTune & \citealp{xue2026fourtune} & 2026 & R & Cone & defect & GL & Mon$\times$Mon & Mq & span & \\
FrameFT & \citealp{adepu2026fine} & 2026 & R & Lin & regular & none & $\Pi$ & M1 & idemp. & \\
$\bullet$\,FRoD & \citealp{wan2025frod} & 2025 & R & Lin & regular & none & O$\times$O & M1 & idemp. & \checkmark\\
GeoRA & \citealp{zhang2026geora} & 2026 & R & Cone$^{*}$ & regular & GL & Mon$\times$Mon & M1 & span & \\
GiVA & \citealp{gangwar2026giva} & 2026 & R & Cone & apex & scalar & O$\times$O & M1 & span & \\
GoRA & \citealp{he2025gora} & 2025 & R & Cone & defect & GL & GL$\times$GL & M1 & span & \\
$\bullet$\,GraLoRA & \citealp{jung2025gralora} & 2025 & R & Cone & apex & GL & $\Pi$ & M1 & span & \checkmark\\
HetLoRA & \citealp{cho2024heterogeneous} & 2024 & R & Cone & apex & GL & GL$\times$GL & M1 & span & \\
Hybrid-LoRA & \citealp{zhang2026hybrid} & 2026 & R & Cone & apex & GL & GL$\times$GL & M1 & span & \\
IC-LoRA & \citealp{huang2024context} & 2024 & R & Cone & apex & GL & GL$\times$GL & M1 & span & \\
IncreLoRA & \citealp{zhang2023increlora} & 2023 & R & Cone & apex & nonlin. & GL$\times$GL & M1 & span & \\
$\bullet$\,InfLoRA & \citealp{liang2024inflora} & 2024 & R & Lin & regular & none & GL$\times$CO & M1 & sched. & \\
Intrinsic SAID & \citealp{aghajanyan2020intrinsic} & 2020 & R & Cone & apex & scalar & $\Pi$ & M1 & span & \\
IR-QLoRA & \citealp{qin2024accurate} & 2024 & R & Cone & defect & GL & Mon$\times$Mon & Mq & span & \\
KAdaptation & \citealp{he2022parameter} & 2022 & R & Cone & apex & GL & $\Pi$ & M1 & span & \\
$\bullet$\,KaSA & \citealp{wang2024kasa} & 2024 & R & Cone & defect & nonlin. & O$\times$O & M1 & span & \checkmark\\
KeepLoRA & \citealp{luo2026keeplora} & 2026 & R & Lin & regular & none & O$\times$O & M1 & sched. & \\
KRAdapter & \citealp{albert2025higher} & 2025 & R & Cone & apex & torus & $\Pi$ & M1 & span & \\
$\bullet$\,KronA & \citealp{edalati2022krona} & 2022 & R & Cone & apex & scalar & $\Pi$ & M1 & span & \\
LaMDA & \citealp{azizi2024lamda} & 2024 & R & Lin & regular & GL & O$\times$O & M1 & idemp. & \\
LoCA & \citealp{du2025loca} & 2025 & R & Lin & regular & none & $\Pi$ & M1 & idemp. & \\
LoDA & \citealp{he2026task} & 2026 & R & Lin & regular & none & GL$\times$CO & M1 & sched. & \\
LoFT & \citealp{tastan2025loft} & 2025 & R & Cone & apex & GL & GL$\times$GL & M1 & span & \\
$\bullet$\,LoftQ & \citealp{li2023loftq} & 2023 & R & Cone$^{*}$ & defect & GL & Mon$\times$Mon & Mq & span & \checkmark\\
$\bullet$\,LoHa & \citealp{hyeonwoo2021fedpara} & 2021 & R & Cone & apex & GL$\times$T & Mon$\times$Mon & M1 & span & \checkmark\\
$\bullet$\,LoKr & \citealp{yeh2023navigating} & 2023 & R & Cone & apex & GL & $\Pi$ & M1 & span & \checkmark\\
$\bullet$\,LongLoRA & \citealp{chen2023longlora} & 2023 & R & Cone & apex & GL & Mon$\times$GL & M1 & span & \\
$\bullet$\,LoRA & \citealp{hu2021lora} & 2021 & R & Cone & apex & GL & GL$\times$GL & M1 & span & \checkmark\\
LoRA Dropout & \citealp{lin2024lora} & 2024 & R & Cone & apex & GL & GL$\times$GL & M1 & span & \\
LoRA Without Regret & \citealp{schulman2025lora} & 2025 & R & Cone & apex & GL & GL$\times$GL & M1 & span & \checkmark\\
$\bullet$\,LoRA+ & \citealp{hayou2024lora} & 2024 & R & Cone & apex & GL & GL$\times$GL & M1 & span & \checkmark\\
LoRA-DA & \citealp{zhang2025lora} & 2025 & R & Cone & defect & GL & GL$\times$GL & M1 & span & \\
LoRA-Dash & \citealp{si2024unleashing} & 2024 & R & Cone & apex & GL & O$\times$O & M1 & span & \\
LoRA-drop & \citealp{zhou2024lora} & 2024 & R & Cone & apex & GL & GL$\times$GL & M1 & span & \\
$\bullet$\,LoRA-FA & \citealp{zhang2023lora} & 2023 & R & Lin & regular & none & frame & M1 & idemp. & \checkmark\\
$\bullet$\,LoRA-GA & \citealp{wang2024lorab} & 2024 & R & Cone$^{*}$ & regular & GL & O$\times$O & M1 & span & \checkmark\\
LoRA-Mini & \citealp{singh2024lora} & 2024 & R & Cone & defect & GL & frame & M1 & span & \\
LoRA-Null & \citealp{tang2025put} & 2025 & R & Cone$^{*}$ & regular & GL & GL$\times$CO & M1 & span & \\
LoRA-One & \citealp{zhang2025lorab} & 2025 & R & Cone & defect & GL & GL$\times$GL & M1 & span & \\
$\bullet$\,LoRA-Pro & \citealp{wang2024lorac} & 2024 & R & Cone & apex & GL & GL$\times$GL & M1 & span & \\
$\bullet$\,LoRA-RITE & \citealp{yen2024lora} & 2024 & R & Cone & apex & GL & GL$\times$GL & M1 & span & \\
LoRA-RLPO / LoRA-RLMO & \citealp{zhang2026geometry} & 2026 & R & Cone & apex & GL & GL$\times$GL & M1 & span & \\
LoRA-S (AdamS, LRACS) & \citealp{zheng2025lora} & 2025 & R & Cone & apex & GL & GL$\times$GL & M1 & span & \\
$\bullet$\,LoRA-SB & \citealp{ponkshe2024initialization} & 2024 & R & Lin & defect & none & Mon$\times$Mon & M1 & idemp. & \\
$\bullet$\,LoRA-XS & \citealp{baazy2024lora} & 2024 & R & Lin & regular & none & O$\times$O & M1 & idemp. & \\
$\bullet$\,LoRETTA & \citealp{yang2024loretta} & 2024 & R & Cone$^{*}$ & regular & GL & $\Pi$ & M1 & span & \\
LoRI & \citealp{zhang2025lori} & 2025 & R & Lin & regular & none & frame & M1 & idemp. & \\
LoRTA & \citealp{hounie2024lorta} & 2024 & R & Cone & apex & torus & $\Pi$ & M1 & span & \\
LoTR & \citealp{bershatsky2024lotr} & 2024 & R & Cone & apex & GL & GL$\times$GL & M1 & span & \\
LQ-LoRA & \citealp{guo2023lq} & 2023 & R & Cone$^{*}$ & defect & GL & Mon$\times$Mon & Mq & span & \\
MatryoshkaLoRA & \citealp{modoranu2026matryoshkalora} & 2026 & R & Cone & apex & GL & GL$\times$GL & M1 & span & \\
MELoRA & \citealp{ren2024melora} & 2024 & R & Cone & apex & GL & $\Pi$ & M1 & span & \\
MetaTT & \citealp{lopezpiqueres2025metatt} & 2025 & R & Cone & apex & GL & GL$\times$GL & M1 & span & \\
$\bullet$\,MiCA & \citealp{rudiger2026mica} & 2026 & R & Lin & regular & none & O$\times$O & M1 & idemp. & \checkmark\\
$\bullet$\,MiLoRA & \citealp{wang2024milora} & 2024 & R & Cone$^{*}$ & regular & GL & O$\times$O & M1 & span & \\
$\bullet$\,MiSS (formerly Bone) & \citealp{kang2024miss} & 2024 & R & Lin & regular & none & $\Pi$ & M1 & idemp. & \checkmark\\
ModuLoRA & \citealp{yin2023modulora} & 2023 & R & Cone & defect & GL & Mon$\times$Mon & Mq & span & \\
MonteCLoRA & \citealp{sengupta2024robust} & 2024 & R & Cone & apex & GL & GL$\times$GL & M1 & span & \checkmark\\
$\bullet$\,MoRA & \citealp{jiang2024mora} & 2024 & R & Lin & regular & none & $\Pi$ & M1 & idemp. & \\
MoRe & \citealp{tan2024more} & 2024 & R & Cone & apex & torus & $\Pi$ & M1 & span & \\
MoS & \citealp{wang2024mos} & 2024 & R & Cone & apex & scalar & $\Pi$ & M1 & span & \\
MoSA & \citealp{wang2026mosa} & 2026 & R & Lin & regular & none & frame & M1 & idemp. & \\
MoSLoRA & \citealp{wu2024mixtureb} & 2024 & R & Cone & apex & GL & GL$\times$GL & M1 & span & \\
$\bullet$\,NOLA & \citealp{koohpayegani2023nola} & 2023 & R & Cone & apex & scalar & frame & M1 & span & \\
$\bullet$\,O-LoRA & \citealp{wang2023orthogonal} & 2023 & R & Cone & apex & GL & GL$\times$GL & M1 & sched. & \\
$\bullet$\,OLoRA & \citealp{buyukakyuz2024olora} & 2024 & R & Cone$^{*}$ & regular & GL & $\Pi$ & M1 & span & \checkmark\\
Orthogonal Adaptation & \citealp{po2023orthogonal} & 2023 & R & Lin & regular & none & frame & M1 & idemp. & \\
$\bullet$\,OSF & \citealp{nayak2025sculpting} & 2025 & R & Lin & regular & nonlin. & O$\times$O & M1 & idemp. & \checkmark\\
$\bullet$\,PEANuT & \citealp{zhong2024peanut} & 2024 & R & Cone & apex & nonlin. & Mon$\times$GL & M1 & span & \checkmark\\
$\bullet$\,PiSSA & \citealp{meng2024pissa} & 2024 & R & Cone$^{*}$ & regular & GL & O$\times$O & M1 & span & \checkmark\\
PLoP & \citealp{hayou2025plop} & 2025 & R & Cone & apex & GL & GL$\times$GL & M1 & span & \\
PoLAR & \citealp{lion2025polar} & 2025 & R & Cone & apex & GL & GL$\times$GL & M1 & span & \\
$\bullet$\,PRISM & \citealp{wang2026prism} & 2026 & R & Cone$^{*}$ & regular & GL & O$\times$O & M1 & span & \\
PRoLoRA & \citealp{wang2024prolora} & 2024 & R & Cone & apex & GL & $\Pi$ & M1 & span & \\
$\bullet$\,PVeRA & \citealp{fillioux2025pvera} & 2025 & R & Cone & apex & scalar & frame & M1 & span & \checkmark\\
$\bullet$\,QA-LoRA & \citealp{xu2023qa} & 2023 & R & Cone & defect & GL & $\Pi$ & M1 & span & \checkmark\\
QeRL & \citealp{huang2025qerl} & 2025 & R & Cone & defect & GL & Mon$\times$Mon & Mq & span & \\
$\bullet$\,QLoRA & \citealp{dettmers2023qlora} & 2023 & R & Cone & defect & GL & Mon$\times$Mon & Mq & span & \checkmark\\
QuanTA & \citealp{chen2024quanta} & 2024 & R & Cone$^{*}$ & regular & GL & $\Pi$ & M1 & span & \\
Quantum-PEFT & \citealp{koikeakino2025quantum} & 2025 & R & Cone & apex & nonlin. & $\Pi$ & M1 & span & \\
$\bullet$\,RandLoRA & \citealp{albert2025randlora} & 2025 & R & Cone & apex & torus & frame & M1 & span & \checkmark\\
RaSA & \citealp{he2025rasa} & 2025 & R & Cone & apex & GL & GL$\times$GL & M1 & span & \\
$\bullet$\,Riemannian LoRA & \citealp{zhang2024riemannian} & 2024 & R & Cone & apex & GL & GL$\times$GL & M1 & span & \checkmark\\
Riemannion & \citealp{bogachev2025lora} & 2025 & R & Cone$^{*}$ & regular & none & O$\times$O & M1 & span & \\
$\bullet$\,RoSA & \citealp{nikdan2024rosa} & 2024 & R & Cone & apex & GL & Mon$\times$Mon & M1 & span & \\
$\bullet$\,rsLoRA & \citealp{kalajdzievski2023rank} & 2023 & R & Cone & apex & GL & GL$\times$GL & M1 & span & \checkmark\\
SD-LoRA & \citealp{wu2025sd} & 2025 & R & Cone & defect & GL & GL$\times$GL & M1 & sched. & \\
SEAL & \citealp{zweiger2025self} & 2025 & R & Cone & apex & GL & GL$\times$GL & M1 & span & \\
ShareLoRA & \citealp{song2024sharelora} & 2024 & R & Cone & apex & GL & GL$\times$GL & M1 & span & \\
Sine-LoRA & \citealp{ji2024efficient} & 2024 & R & Cone & apex & GL & Mon$\times$Mon & M1 & span & \\
SingLoRA & \citealp{bensaid2025singlora} & 2025 & R & Cone & defect & GL & $\Pi$ & M1 & span & \\
SLAO (Merge before Forget) & \citealp{qiao2025merge} & 2025 & R & Cone & defect & GL & GL$\times$GL & M1 & span & \\
SLTrain & \citealp{han2024sltrain} & 2024 & R & Cone & defect & GL & Mon$\times$Mon & M1 & span & \\
SoRA & \citealp{ding2023sparse} & 2023 & R & Cone & apex & nonlin. & GL$\times$GL & M1 & span & \\
SparseLoRA & \citealp{khaki2025sparselora} & 2025 & R & Cone & apex & GL & GL$\times$GL & M1 & span & \\
Spectral Adapter\textsuperscript{A} & \citealp{zhang2024spectral} & 2024 & R & Cone$^{*}$ & regular & GL & O$\times$O & M1 & span & \\
SRR (Structured Residual Reconstruction) & \citealp{cho2026preserve} & 2026 & R & Cone$^{*}$ & defect & GL & Mon$\times$Mon & Mq & span & \\
SSH & \citealp{shen2025ssh} & 2025 & R & Lin & regular & none & $\Pi$ & M1 & idemp. & \\
Stable-LoRA & \citealp{wu2025stable} & 2025 & R & Cone & apex & GL & GL$\times$GL & M1 & span & \\
StructLoRA & \citealp{xiao2026not} & 2026 & R & Lin & defect & nonlin. & GL$\times$GL & M1 & idemp. & \\
$\bullet$\,SVFT & \citealp{lingam2024svft} & 2024 & R & Lin & regular & none & O$\times$O & M1 & idemp. & \\
TeRA & \citealp{gu2025tera} & 2025 & R & Cone & apex & torus & $\Pi$ & M1 & span & \\
Tied-LoRA & \citealp{renduchintala2023tied} & 2023 & R & Cone & apex & torus & Mon$\times$GL & M1 & span & \\
Tina & \citealp{wang2025tina} & 2025 & R & Cone & apex & GL & GL$\times$GL & M1 & span & \checkmark\\
$\bullet$\,TinyLoRA & \citealp{morris2026learning} & 2026 & R & Lin & regular & none & O$\times$O & M1 & idemp. & \checkmark\\
Trans-LoRA & \citealp{wang2024trans} & 2024 & R & Cone & apex & GL & GL$\times$GL & M1 & span & \\
TT-LoRA & \citealp{anjum2024tensor} & 2024 & R & Cone & defect & GL & $\Pi$ & M1 & span & \\
$\bullet$\,Uni-LoRA & \citealp{li2025uni} & 2025 & R & Cone & defect & none & frame & M1 & span & \checkmark\\
$\bullet$\,VB-LoRA & \citealp{li2024vb} & 2024 & R & Cone & defect & none & $\Pi$ & M1 & span & \checkmark\\
VeLoRA & \citealp{miles2024velora} & 2024 & R & Cone & apex & GL & GL$\times$GL & M1 & span & \checkmark\\
$\bullet$\,VeRA & \citealp{kopiczko2023vera} & 2023 & R & Cone & apex & scalar & frame & M1 & span & \checkmark\\
WASI & \citealp{nguyen2025efficient} & 2025 & R & Cone$^{*}$ & defect & GL & O$\times$O & M1 & idemp. & \\
$\bullet$\,WaveFT & \citealp{bilican2025exploring} & 2025 & R & Lin & regular & none & $\Pi$ & M1 & idemp. & \checkmark\\
\multicolumn{11}{@{}l}{\sffamily\bfseries\color{tide}R, action}\\[1pt]
$\bullet$\,(IA)\textsuperscript{3} & \citealp{liu2022few} & 2022 & R & Lin & regular & none & Mon$\times$GL & M1$\to$ & idemp. & \checkmark\\
BiDoRA & \citealp{qin2024bidora} & 2024 & R & Sat & apex & GL & Mon$\times$GL & M1 & span & \\
$\bullet$\,BOFT & \citealp{liu2023parameter} & 2023 & R & Sat & regular & nonlin. & $\Pi$ & M1 & group & \checkmark\\
$\bullet$\,DoRA & \citealp{liu2024dora} & 2024 & R & Sat & apex & GL & Mon$\times$GL & M1 & span & \checkmark\\
EfficientQAT & \citealp{chen2024efficientqat} & 2024 & R & Lin & defect & none & $\Pi$ & M1 & idemp. & \\
ETHER & \citealp{bini2024ether} & 2024 & R & Orb & unpt. & torus & GL$\times$CO & M1 & group & \\
ETHER+ & \citealp{bini2024ether} & 2024 & R & Cone & apex & torus & O$\times$O & M1 & group & \\
$\bullet$\,GOFT & \citealp{ma2024parameter} & 2024 & R & Orb & regular & none & $\Pi$ & M1 & group & \\
$\bullet$\,HRA & \citealp{yuan2024bridging} & 2024 & R & Meet & apex & nonlin. & GL$\times$CO & M1 & group & \checkmark\\
L4Q & \citealp{jeon2024l4q} & 2024 & R & Lin & defect & nonlin. & $\Pi$ & M1 & span & \\
$\bullet$\,LoCO & \citealp{nguyen2026loco} & 2026 & R & Meet & apex & GL & GL$\times$CO & M1 & group & \\
$\bullet$\,LOFT & \citealp{zhao2026loft} & 2026 & R & Orb & regular & none & O$\times$O & M1 & idemp. & \\
LoRA-X & \citealp{farhadzadeh2025lora} & 2025 & R & Lin & regular & none & O$\times$O & M1 & idemp. & \\
LoRAPrune & \citealp{zhang2023loraprune} & 2023 & R & Sat & apex & GL & Mon$\times$GL & M1 & span & \\
Masking & \citealp{zhao2020masking} & 2020 & R & Orb & defect & torus & Mon$\times$Mon & M1 & idemp. & \\
$\bullet$\,OFT & \citealp{qiu2023controlling} & 2023 & R & Orb & regular & none & GL$\times$CO & M1 & idemp. & \checkmark\\
$\bullet$\,OFTv2 & \citealp{qiu2025orthogonal} & 2025 & R & Orb & regular & none & GL$\times$CO & M1 & group & \checkmark\\
PaRa & \citealp{chen2024para} & 2024 & R & Orb & unpt. & GL & GL$\times$CO & M1 & group & \checkmark\\
$\bullet$\,PEQA & \citealp{kim2023memory} & 2023 & R & Lin & defect & none & Mon$\times$Mon & M1 & idemp. & \\
$\bullet$\,Piggyback & \citealp{mallya2018piggyback} & 2018 & R & Orb & regular & torus & Mon$\times$Mon & M1 & idemp. & \\
$\bullet$\,PSOFT & \citealp{wu2025efficient} & 2025 & R & Sat & apex & scalar & O$\times$O & M1 & group & \checkmark\\
$\bullet$\,qGOFT & \citealp{ma2024parameter} & 2024 & R & Orb & regular & torus & $\Pi$ & M1 & group & \\
QZO & \citealp{shang2025fine} & 2025 & R & Lin & defect & none & $\Pi$ & M1 & idemp. & \\
$\bullet$\,RED & \citealp{wu2024advancing} & 2024 & R & Lin & regular & none & Mon$\times$GL & M1$\to$ & idemp. & \\
RepAdapter & \citealp{luo2023efficient} & 2023 & R & Cone & apex & GL & $\Pi$ & M1$\to$ & group & \\
Residual Adapters & \citealp{rebuffi2017learning} & 2017 & R & Lin & regular & torus & GL$\times$GL & M1$\to$ & idemp. & \\
$\bullet$\,RoAd & \citealp{liao20243} & 2024 & R & Orb & regular & none & Mon$\times$GL & M1 & idemp. & \checkmark\\
Spectral Adapter\textsuperscript{R} & \citealp{zhang2024spectral} & 2024 & R & Orb & regular & none & O$\times$O & M1 & idemp. & \\
$\bullet$\,SSF & \citealp{lian2022scaling} & 2022 & R & Lin & regular & none & Mon$\times$GL & M1$\to$ & idemp. & \\
SVDiff & \citealp{han2023svdiff} & 2023 & R & Lin & regular & none & O$\times$O & M1 & idemp. & \\
SVF (Transformer\textsuperscript{2}) & \citealp{sun2025transformer} & 2025 & R & Lin & regular & none & O$\times$O & M1 & idemp. & \\
\multicolumn{11}{@{}l}{\sffamily\bfseries\color{tide}R, Hadamard}\\[1pt]
FLoRA & \citealp{wen2023batched} & 2023 & R & Cone & apex & GL & Mon$\times$Mon & M1 & group & \\
$\bullet$\,GLoRA & \citealp{chavan2023one} & 2023 & R & Cone & apex & GL & Mon$\times$Mon & M1 & span & \checkmark\\
$\bullet$\,HiRA & \citealp{huang2025hira} & 2025 & R & Cone & apex & GL & Mon$\times$Mon & M1 & span & \checkmark\\
LieRA & \citealp{si2025generalized} & 2025 & R & Cone & apex & GL & Mon$\times$Mon & M1 & span & \\
Low-Rank Ternary Adaptation & \citealp{manolache2026low} & 2026 & R & Orb & regular & nonlin. & $\Pi$ & M1 & idemp. & \\
SMoA & \citealp{liu2026smoa} & 2026 & R & Cone & apex & GL & Mon$\times$Mon & M1 & span & \\
\multicolumn{11}{@{}l}{\sffamily\bfseries\color{tide}R, selective}\\[1pt]
$\bullet$\,BEFT & \citealp{huang2025beft} & 2025 & R & Lin & regular & none & GL$\times$GL & M1 & idemp. & \checkmark\\
$\bullet$\,BitFit & \citealp{benzaken2021bitfit} & 2021 & R & Lin & regular & none & GL$\times$GL & M1 & idemp. & \\
$\bullet$\,Child-Tuning & \citealp{xu2021raise} & 2021 & R & Lin & regular & none & Mon$\times$Mon & M1 & idemp. & \\
DiaBlo & \citealp{gurses2025diablo} & 2025 & R & Lin & regular & none & $\Pi$ & M1 & idemp. & \\
$\bullet$\,Diff Pruning & \citealp{guo2020parameter} & 2020 & R & Cone & apex & none & Mon$\times$Mon & M1 & span & \\
DiffFit & \citealp{xie2023difffit} & 2023 & R & Lin & regular & none & Mon$\times$GL & M1$\to$ & idemp. & \\
ESFT & \citealp{wang2024let} & 2024 & R & Lin & regular & none & GL$\times$GL & M1 & idemp. & \\
$\bullet$\,FISH Mask & \citealp{sung2021training} & 2021 & R & Lin & regular & none & Mon$\times$Mon & M1 & idemp. & \\
GPS & \citealp{zhang2023gradient} & 2023 & R & Lin & regular & none & Mon$\times$Mon & M1 & idemp. & \\
KappaTune & \citealp{ludwig2025condition} & 2025 & R & Lin & regular & none & O$\times$O & M1 & idemp. & \checkmark\\
$\bullet$\,Linear probing & \citealp{donahue2013decaf} & 2013 & R & Lin & regular & none & GL$\times$GL & M1 & idemp. & \\
$\bullet$\,LN Tuning & \citealp{zhao2023tuning} & 2023 & R & Lin & regular & none & Mon$\times$GL & M1$\to$ & idemp. & \checkmark\\
LoFiT & \citealp{yin2024lofit} & 2024 & R & Lin & regular & none & GL$\times$GL & M1$\to$ & idemp. & \\
$\bullet$\,LT-SFT & \citealp{ansell2021composable} & 2021 & R & Lin & regular & none & Mon$\times$Mon & M1 & idemp. & \\
MDPD & \citealp{zhang2026memory} & 2026 & R & Lin & regular & none & Mon$\times$GL & M1$\to$ & idemp. & \\
PaFi & \citealp{liao2023parameter} & 2023 & R & Lin & regular & none & Mon$\times$Mon & M1 & idemp. & \\
SAM (sparse) & \citealp{fu2022effectiveness} & 2022 & R & Lin & regular & none & Mon$\times$Mon & M1 & idemp. & \\
$\bullet$\,SHiRA & \citealp{bhardwaj2024sparse} & 2024 & R & Lin & regular & none & Mon$\times$Mon & M1 & idemp. & \checkmark\\
SIFT & \citealp{song2023sparse} & 2023 & R & Lin & regular & none & Mon$\times$Mon & M1 & idemp. & \\
SMT & \citealp{he2024smt} & 2024 & R & Lin & regular & none & $\Pi$ & M1 & idemp. & \\
Spectrum & \citealp{hartford2024spectrum} & 2024 & R & Lin & regular & none & O$\times$O & M1 & idemp. & \\
SpIEL & \citealp{ansell2024scaling} & 2024 & R & Cone & apex & none & Mon$\times$Mon & M1 & span & \\
SPT & \citealp{he2023sensitivity} & 2023 & R & Cone & apex & GL & Mon$\times$Mon & M1 & span & \\
$\bullet$\,Super-Tuning (Super / Supra) & \citealp{ilin2026super} & 2026 & R & Lin & regular & none & Mon$\times$Mon & M1 & idemp. & \checkmark\\
$\bullet$\,Surgical FT & \citealp{lee2022surgical} & 2022 & R & Lin & regular & none & GL$\times$GL & M1 & idemp. & \\
S\textsuperscript{2}FT & \citealp{yang2024s2ft} & 2024 & R & Lin & regular & none & Mon$\times$GL & M1 & idemp. & \\
TinyTL & \citealp{cai2020tinytl} & 2020 & R & Lin & defect & none & Mon$\times$Mon & M1 & idemp. & \\
$\bullet$\,Trainable Tokens & \citealp{team2025trainable} & 2025 & R & Lin & regular & none & Mon$\times$GL & M1 & idemp. & \checkmark\\
\multicolumn{11}{@{}l}{\sffamily\bfseries\color{tide}E (pointed ext.)}\\[1pt]
AdaMix & \citealp{wang2022adamix} & 2022 & E & Fun & neutral & none & GL$\times$GL & M$\infty$ & -- & \\
$\bullet$\,Adapter (Houlsby) & \citealp{houlsby2019parameter} & 2019 & E & Fun & neutral & none & GL$\times$GL & M$\infty$ & -- & \\
Adapter+ & \citealp{steitz2024adapters} & 2024 & E & Fun & neutral & torus & GL$\times$GL & M$\infty$ & -- & \\
AdapterDrop & \citealp{ruckle2020adapterdrop} & 2020 & E & Fun & neutral & torus & GL$\times$GL & M$\infty$ & -- & \\
AdaptFormer & \citealp{chen2022adaptformer} & 2022 & E & Fun & neutral & torus & GL$\times$GL & M$\infty$ & -- & \\
AIM & \citealp{yang2023aim} & 2023 & E & Fun & neutral & none & GL$\times$GL & M$\infty$ & -- & \\
AP (Activation Prompt) & \citealp{zhang2023visual} & 2023 & E & Lin & neutral & none & GL$\times$GL & M$\infty$ & -- & \\
CLIP-Adapter & \citealp{gao2021clip} & 2021 & E & Fun & defect & torus & GL$\times$GL & M$\infty$ & -- & \\
$\bullet$\,Compacter & \citealp{mahabadi2021compacter} & 2021 & E & Fun & neutral & torus & $\Pi$ & M$\infty$ & -- & \\
$\bullet$\,ControlNet & \citealp{zhang2023adding} & 2023 & E & Fun & neutral & GL & GL$\times$GL & M$\infty$ & -- & \\
Convpass & \citealp{jie2022convolutional} & 2022 & E & Fun & neutral & none & GL$\times$GL & M$\infty$ & -- & \\
$\bullet$\,Custom Diffusion & \citealp{kumari2022multi} & 2022 & E & Lin & neutral & none & GL$\times$GL & M1 & idemp. & \\
Initial State Tuning & \citealp{galim2024parameter} & 2024 & E & Lin & neutral & none & GL$\times$GL & M$\infty$ & idemp. & \\
IP-Adapter & \citealp{ye2023ip} & 2023 & E & Fun & defect & torus & GL$\times$GL & M$\infty$ & -- & \\
K-Adapter & \citealp{wang2020k} & 2020 & E & Fun & neutral & GL & GL$\times$GL & M$\infty$ & -- & \\
$\bullet$\,LLaMA-Adapter & \citealp{zhang2023llama} & 2023 & E & Fun & neutral & none & GL$\times$GL & M$\infty$ & -- & \checkmark\\
LLaMA-Adapter V2 & \citealp{gao2023llama} & 2023 & E & Fun & neutral & torus & Mon$\times$GL & M$\infty$ & -- & \\
$\bullet$\,LoReFT & \citealp{wu2024reft} & 2024 & E & Fun & defect & GL & GL$\times$GL & M$\infty$ & -- & \\
MAD-X & \citealp{pfeiffer2020mad} & 2020 & E & Fun & neutral & torus & $\Pi$ & M$\infty$ & -- & \\
Mona & \citealp{yin20245} & 2024 & E & Fun & defect & torus & Mon$\times$Mon & M$\infty$ & -- & \\
$\bullet$\,Parallel Adapter & \citealp{he2021unified} & 2021 & E & Fun & neutral & torus & GL$\times$GL & M$\infty$ & -- & \\
$\bullet$\,Pfeiffer Adapter & \citealp{pfeiffer2020adapterfusion} & 2020 & E & Fun & neutral & torus & GL$\times$GL & M$\infty$ & -- & \\
PHM-Adapter & \citealp{mahabadi2021compacter} & 2021 & E & Fun & neutral & GL & $\Pi$ & M$\infty$ & -- & \\
PrefixMemory-Tuning & \citealp{wang2025prefixmemory} & 2025 & E & Fun & neutral & none & GL$\times$GL & M$\infty$ & -- & \\
Res-Tuning & \citealp{jiang2023res} & 2023 & E & Fun & neutral & none & GL$\times$GL & M$\infty$ & -- & \\
$\bullet$\,ShadowPEFT & \citealp{li2026shadowpeft} & 2026 & E & Fun & neutral & GL & Mon$\times$Mon & M$\infty$ & -- & \checkmark\\
Side-Tuning & \citealp{zhang2019side} & 2019 & E & Fun & neutral & torus & GL$\times$GL & M$\infty$ & -- & \\
ST-Adapter & \citealp{pan2022st} & 2022 & E & Fun & neutral & torus & GL$\times$GL & M$\infty$ & -- & \\
State-offset Tuning & \citealp{kang2025state} & 2025 & E & Lin & neutral & none & GL$\times$GL & M$\infty$ & idemp. & \\
T2I-Adapter & \citealp{mou2023t2i} & 2023 & E & Fun & defect & torus & GL$\times$GL & M$\infty$ & -- & \\
$\bullet$\,Textual Inversion & \citealp{gal2022image} & 2022 & E & Lin & neutral & none & Mon$\times$GL & M1 & idemp. & \\
Tip-Adapter & \citealp{zhang2022tip} & 2022 & E & Fun & defect & none & O$\times$O & M$\infty$ & -- & \\
Visual Prompting & \citealp{bahng2022exploring} & 2022 & E & Lin & defect & none & $\Pi$ & M1$\to$ & idemp. & \\
\multicolumn{11}{@{}l}{\sffamily\bfseries\color{tide}E$^{0}$ (unpointed)}\\[1pt]
AutoPEFT & \citealp{zhou2023autopeft} & 2023 & E$^{0}$ & Fun & unpt. & torus & GL$\times$GL & M$\infty$ & -- & \\
$\bullet$\,AutoPrompt & \citealp{shin2020autoprompt} & 2020 & E$^{0}$ & Fun & unpt. & none & GL$\times$GL & M$\infty$ & -- & \\
BBT & \citealp{sun2022black} & 2022 & E$^{0}$ & Fun & unpt. & none & frame & M$\infty$ & -- & \\
$\bullet$\,Cartridges & \citealp{eyuboglu2025cartridges} & 2025 & E$^{0}$ & Fun & unpt. & none & GL$\times$GL & M$\infty$ & -- & \checkmark\\
$\bullet$\,CoOp & \citealp{zhou2021learning} & 2021 & E$^{0}$ & Fun & unpt. & none & GL$\times$GL & M1$\to$ & -- & \\
$\bullet$\,CPT & \citealp{blau2024context} & 2024 & E$^{0}$ & Fun & unpt. & none & O$\times$O & M$\infty$ & -- & \checkmark\\
DePT & \citealp{shi2023dept} & 2023 & E$^{0}$ & Fun & unpt. & GL & GL$\times$GL & M$\infty$ & -- & \\
EasyControl & \citealp{zhang2025easycontrol} & 2025 & E$^{0}$ & Fun & unpt. & GL & GL$\times$GL & M$\infty$ & -- & \\
LoRA layer replication (DUS) & \citealp{kim2023solar} & 2023 & E$^{0}$ & Fun & unpt. & GL & GL$\times$GL & M1 & span & \checkmark\\
$\bullet$\,LST & \citealp{sung2022lst} & 2022 & E$^{0}$ & Fun & unpt. & GL & GL$\times$GL & M$\infty$ & -- & \\
$\bullet$\,MAM Adapter & \citealp{he2021unified} & 2021 & E$^{0}$ & Fun & unpt. & GL & GL$\times$GL & M$\infty$ & -- & \\
MaPLe & \citealp{khattak2022maple} & 2022 & E$^{0}$ & Fun & unpt. & nonlin. & GL$\times$GL & M$\infty$ & -- & \\
MEFT & \citealp{liao2023make} & 2023 & E$^{0}$ & Fun & unpt. & none & GL$\times$GL & M$\infty$ & -- & \\
$\bullet$\,MPT & \citealp{wang2023multitask} & 2023 & E$^{0}$ & Fun & unpt. & torus & GL$\times$GL & M$\infty$ & -- & \checkmark\\
NOAH & \citealp{zhang2022neural} & 2022 & E$^{0}$ & Fun & unpt. & GL & GL$\times$GL & M$\infty$ & -- & \\
OminiControl & \citealp{tan2024ominicontrol} & 2024 & E$^{0}$ & Fun & unpt. & GL & GL$\times$GL & M$\infty$ & -- & \\
$\bullet$\,P-Tuning & \citealp{liu2021gpt} & 2021 & E$^{0}$ & Fun & unpt. & nonlin. & GL$\times$GL & M$\infty$ & -- & \checkmark\\
$\bullet$\,P-Tuning v2 & \citealp{liu2021p} & 2021 & E$^{0}$ & Fun & unpt. & none & GL$\times$GL & M$\infty$ & -- & \checkmark\\
PPT & \citealp{gu2021ppt} & 2021 & E$^{0}$ & Fun & unpt. & none & GL$\times$GL & M$\infty$ & -- & \\
$\bullet$\,Prefix-Tuning & \citealp{li2021prefix} & 2021 & E$^{0}$ & Fun & unpt. & nonlin. & GL$\times$GL & M$\infty$ & -- & \checkmark\\
$\bullet$\,Prompt Tuning & \citealp{lester2021power} & 2021 & E$^{0}$ & Fun & unpt. & none & GL$\times$GL & M$\infty$ & -- & \checkmark\\
Residual Prompt Tuning & \citealp{razdaibiedina2023residual} & 2023 & E$^{0}$ & Fun & unpt. & nonlin. & GL$\times$GL & M$\infty$ & -- & \\
SPoT & \citealp{vu2021spot} & 2021 & E$^{0}$ & Fun & unpt. & none & GL$\times$GL & M$\infty$ & -- & \\
$\bullet$\,UniPELT & \citealp{mao2021unipelt} & 2021 & E$^{0}$ & Fun & unpt. & GL & GL$\times$GL & M$\infty$ & -- & \\
$\bullet$\,VPT & \citealp{jia2022visual} & 2022 & E$^{0}$ & Fun & unpt. & torus & GL$\times$GL & M$\infty$ & -- & \\
WARP & \citealp{hambardzumyan2021warp} & 2021 & E$^{0}$ & Fun & unpt. & none & GL$\times$GL & M$\infty$ & -- & \\
XPrompt & \citealp{ma2022xprompt} & 2022 & E$^{0}$ & Fun & unpt. & none & $\Pi$ & M$\infty$ & -- & \\
\multicolumn{11}{@{}l}{\sffamily\bfseries\color{tide}B (backward)}\\[1pt]
AdaLomo & \citealp{lv2023adalomo} & 2023 & B & Lin & regular & none & GL$\times$GL & na & idemp. & \\
Adam-mini & \citealp{zhang2024adam} & 2024 & B & Lin & regular & none & GL$\times$GL & na & idemp. & \\
AGZO & \citealp{lin2026agzo} & 2026 & B & Lin & regular & none & frame & -- & sched. & \\
APOLLO & \citealp{zhu2024apollo} & 2024 & B & Lin & regular & none & GL$\times$GL & na & idemp. & \\
BAdam & \citealp{luo2024badam} & 2024 & B & Lin & regular & none & GL$\times$GL & na & sched. & \\
COLA & \citealp{xia2024chain} & 2024 & B & Cone & apex & GL & GL$\times$GL & M1 & sched. & \\
Delta-LoRA & \citealp{zi2023delta} & 2023 & B & Cone & apex & GL & GL$\times$GL & M1 & sched. & \\
Fira & \citealp{chen2024fira} & 2024 & B & Lin & regular & none & GL$\times$GL & na & idemp. & \\
$\bullet$\,Flora & \citealp{hao2024flora} & 2024 & B & Lin & regular & none & frame & na & sched. & \\
FLoRA (federated) & \citealp{wang2024flora} & 2024 & B & Cone & apex & GL & GL$\times$GL & M1 & sched. & \\
FZOO & \citealp{dang2025fzoo} & 2025 & B & Lin & regular & none & Mon$\times$Mon & -- & sched. & \\
$\bullet$\,GaLore & \citealp{zhao2024galore} & 2024 & B & Lin & regular & none & O$\times$O & na & sched. & \\
GaLore 2 & \citealp{su2025galore} & 2025 & B & Lin & regular & none & O$\times$O & na & sched. & \\
HiFT & \citealp{liu2024hift} & 2024 & B & Lin & regular & none & GL$\times$GL & na & sched. & \\
LIFT & \citealp{liu2025lift} & 2025 & B & Lin & regular & none & Mon$\times$Mon & na & sched. & \\
$\bullet$\,LISA & \citealp{pan2024lisa} & 2024 & B & Lin & regular & none & GL$\times$GL & na & sched. & \\
LOMO & \citealp{lv2023full} & 2023 & B & Lin & regular & none & GL$\times$GL & na & idemp. & \\
LoQT & \citealp{loeschcke2024loqt} & 2024 & B & Lin & defect & none & Mon$\times$Mon & Mq & sched. & \\
LoRA-Pre & \citealp{wang2026taming} & 2026 & B & Lin & regular & none & GL$\times$GL & -- & -- & \\
$\bullet$\,LP-FT & \citealp{kumar2022fine} & 2022 & B & Lin & regular & none & GL$\times$GL & na & sched. & \\
$\bullet$\,MeZO & \citealp{malladi2023fine} & 2023 & B & Lin & regular & none & frame & na & sched. & \\
Online Subspace Descent & \citealp{liang2024memory} & 2024 & B & Lin & regular & none & O$\times$O & na & sched. & \\
PeriodicLoRA & \citealp{meng2024periodiclora} & 2024 & B & Cone & apex & GL & GL$\times$GL & M1 & sched. & \\
POET-X & \citealp{qiu2026poet} & 2026 & B & Orb & regular & none & Mon$\times$Mon & M1 & sched. & \\
Q-GaLore & \citealp{zhang2024q} & 2024 & B & Lin & defect & none & Mon$\times$Mon & na & sched. & \\
$\bullet$\,ReLoRA & \citealp{lialin2023relora} & 2023 & B & Cone & apex & GL & GL$\times$GL & M1 & sched. & \\
ScaLoRA & \citealp{zhang2025scalora} & 2025 & B & Cone$^{*}$ & regular & GL & GL$\times$GL & M1 & sched. & \\
Sparse Memory Finetuning & \citealp{lin2025continual} & 2025 & B & Lin & regular & none & Mon$\times$GL & na & sched. & \\
ZO-Muon & \citealp{lang2026powering} & 2026 & B & Lin & regular & none & frame & -- & sched. & \\
\multicolumn{11}{@{}l}{\sffamily\bfseries\color{tide}P (post hoc)}\\[1pt]
AdaMerging & \citealp{yang2023adamerging} & 2023 & P & Lin & defect & none & GL$\times$GL & M1 & idemp. & \\
AdapterSoup & \citealp{chronopoulou2023adaptersoup} & 2023 & P & Fun & unpt. & -- & -- & M$\infty$ & -- & \\
AdaRank & \citealp{lee2025adarank} & 2025 & P & Cone & defect & torus & O$\times$O & M1 & span & \\
Compress-then-Merge (CtM) & \citealp{he2026compress} & 2026 & P & Lin & regular & none & O$\times$O & M1 & idemp. & \\
Core Space & \citealp{panariello2025accurate} & 2025 & P & Lin & regular & none & O$\times$O & M1 & idemp. & \\
$\bullet$\,DARE & \citealp{yu2023language} & 2023 & P & Lin & regular & none & Mon$\times$Mon & M1 & idemp. & \checkmark\\
DC-Merge & \citealp{zhang2026dc} & 2026 & P & Lin & regular & none & O$\times$O & M1 & idemp. & \\
EFT & \citealp{mitchell2023emulator} & 2023 & P & Fun & neutral & -- & GL$\times$GL & M$\infty$ & -- & \\
Fisher Merging & \citealp{matena2021merging} & 2021 & P & -- & unpt. & -- & Mon$\times$Mon & M1 & -- & \\
GenKnowSub & \citealp{bagherifard2025genknowsub} & 2025 & P & Fun & unpt. & -- & O$\times$O & M$\infty$ & -- & \checkmark\\
GradFix & \citealp{rinaldi2025gradient} & 2025 & P & Lin & regular & none & Mon$\times$Mon & M1 & idemp. & \\
$\bullet$\,Intruder-dimension reduction & \citealp{shuttleworth2024lora} & 2024 & P & Lin & defect & none & O$\times$O & M1 & idemp. & \checkmark\\
$\bullet$\,KnOTS & \citealp{stoica2024model} & 2024 & P & Lin & regular & none & O$\times$O & M1 & idemp. & \\
$\bullet$\,LCM-LoRA & \citealp{luo2023lcm} & 2023 & P & Lin & regular & none & GL$\times$GL & M1 & idemp. & \checkmark\\
$\bullet$\,LoRA Soups (CAT) & \citealp{prabhakar2024lora} & 2024 & P & Lin & regular & none & GL$\times$GL & M1 & idemp. & \\
LoRA Switch / LoRA Composite & \citealp{zhong2024multi} & 2024 & P & Fun & unpt. & -- & -- & M$\infty$ & -- & \\
LoRA-LEGO & \citealp{zhao2024merging} & 2024 & P & -- & unpt. & -- & O$\times$O & M1 & -- & \\
$\bullet$\,LoraHub & \citealp{huang2023lorahub} & 2023 & P & Cone & apex & none & GL$\times$GL & M1 & -- & \\
Mix-of-Show & \citealp{gu2023mix} & 2023 & P & -- & unpt. & -- & GL$\times$GL & M1 & -- & \\
$\bullet$\,Model Soups & \citealp{wortsman2022model} & 2022 & P & -- & unpt. & -- & GL$\times$GL & M1 & -- & \\
OrthoMerge & \citealp{yang2026orthogonal} & 2026 & P & -- & unpt. & -- & GL$\times$CO & M1 & -- & \\
PEM Arithmetic & \citealp{zhang2023composing} & 2023 & P & -- & unpt. & -- & GL$\times$GL & M1 & -- & \checkmark\\
Pico & \citealp{tang2026crowded} & 2026 & P & -- & unpt. & -- & O$\times$O & M1 & -- & \\
$\bullet$\,Proxy-Tuning & \citealp{liu2024tuning} & 2024 & P & Fun & neutral & -- & GL$\times$GL & M$\infty$ & -- & \\
RAIN-Merging & \citealp{huang2025rain} & 2025 & P & Lin & regular & none & GL$\times$CO & M1 & idemp. & \\
RegMean & \citealp{jin2022dataless} & 2022 & P & -- & unpt. & -- & Mon$\times$GL & M1 & -- & \\
$\bullet$\,Task Arithmetic & \citealp{ilharco2022editing} & 2022 & P & Lin & regular & none & GL$\times$GL & M1 & idemp. & \checkmark\\
$\bullet$\,TIES-Merging & \citealp{yadav2023ties} & 2023 & P & Lin & regular & none & Mon$\times$Mon & M1 & idemp. & \checkmark\\
TSV-Merge & \citealp{gargiulo2024task} & 2024 & P & Lin & regular & none & O$\times$O & M1 & idemp. & \\
$\bullet$\,WiSE-FT & \citealp{wortsman2021robust} & 2021 & P & Lin & regular & none & GL$\times$GL & M1 & idemp. & \checkmark\\
ZipLoRA & \citealp{shah2023ziplora} & 2023 & P & Lin & defect & none & Mon$\times$GL & M1 & idemp. & \\
\multicolumn{11}{@{}l}{\sffamily\bfseries\color{tide}F$_T$ (family)}\\[1pt]
$\bullet$\,AdapterFusion & \citealp{pfeiffer2020adapterfusion} & 2020 & F$_T$ & Fun & defect & GL & GL$\times$GL & M$\infty$ & -- & \\
$\bullet$\,aLoRA & \citealp{greenewald2025activated} & 2025 & F$_T$ & Fun & apex & GL & GL$\times$GL & M$\infty$ & -- & \checkmark\\
$\bullet$\,Arrow & \citealp{ostapenko2024modular} & 2024 & F$_T$ & Fun & unpt. & -- & -- & M$\infty$ & -- & \checkmark\\
ATTEMPT & \citealp{asai2022attempt} & 2022 & F$_T$ & Fun & unpt. & nonlin. & Mon$\times$Mon & M$\infty$ & -- & \\
Bayesian-LoRA & \citealp{lin2026bayesian} & 2026 & F$_T$ & Cone & defect & nonlin. & GL$\times$GL & M$\infty$ & -- & \\
BLoB & \citealp{wang2024blob} & 2024 & F$_T$ & Cone & apex & torus & GL$\times$GL & M$\infty$ & -- & \\
CoA-LoRA & \citealp{ye2025fly} & 2025 & F$_T$ & Cone$^{*}$ & defect & GL & Mon$\times$Mon & Mq & -- & \\
CoCoOp & \citealp{zhou2022conditional} & 2022 & F$_T$ & Fun & unpt. & torus & GL$\times$GL & M$\infty$ & -- & \\
CODA-Prompt & \citealp{smith2022coda} & 2022 & F$_T$ & Fun & unpt. & torus & Mon$\times$GL & M$\infty$ & -- & \\
DnD & \citealp{liang2025drag} & 2025 & F$_T$ & Cone & defect & GL & GL$\times$GL & M1 & -- & \\
Doc-to-LoRA (D2L) & \citealp{charakorn2026doc} & 2026 & F$_T$ & Cone & defect & GL & GL$\times$GL & M1 & -- & \\
DualPrompt & \citealp{wang2022dualprompt} & 2022 & F$_T$ & Fun & unpt. & torus & O$\times$O & M$\infty$ & -- & \\
FourierMoE & \citealp{jiang2026fouriermoe} & 2026 & F$_T$ & Fun & neutral & none & $\Pi$ & M$\infty$ & -- & \\
GenerativeAdapter & \citealp{chen2024generative} & 2024 & F$_T$ & Cone & defect & GL & GL$\times$GL & M1 & -- & \\
GOAT & \citealp{fan2025make} & 2025 & F$_T$ & Fun & defect & GL & -- & M$\infty$ & -- & \\
$\bullet$\,HydraLoRA & \citealp{tian2024hydralora} & 2024 & F$_T$ & Fun & neutral & GL & -- & M$\infty$ & -- & \\
$\bullet$\,HyperDreamBooth & \citealp{ruiz2023hyperdreambooth} & 2023 & F$_T$ & Cone & defect & GL & frame & M1 & sched. & \\
$\bullet$\,HyperFormer & \citealp{mahabadi2021parameter} & 2021 & F$_T$ & Fun & defect & GL & GL$\times$GL & M$\infty$ & -- & \\
ICEdit & \citealp{zhang2025context} & 2025 & F$_T$ & Fun & neutral & GL & -- & M$\infty$ & -- & \\
IDPG & \citealp{wu2022idpg} & 2022 & F$_T$ & Fun & unpt. & GL & $\Pi$ & M$\infty$ & -- & \\
In-Place TTT & \citealp{feng2025place} & 2025 & F$_T$ & Fun & neutral & torus & GL$\times$CO & M$\infty$ & -- & \\
K-LoRA & \citealp{ouyang2025k} & 2025 & F$_T$ & -- & unpt. & -- & Mon$\times$Mon & M1 & -- & \\
L2P & \citealp{wang2021learning} & 2021 & F$_T$ & Fun & unpt. & torus & O$\times$O & M$\infty$ & -- & \\
$\bullet$\,Laplace-LoRA & \citealp{yang2023bayesian} & 2023 & F$_T$ & Cone & apex & GL & GL$\times$GL & M$\infty$ & -- & \\
LD-MoLE & \citealp{zhuang2025ld} & 2025 & F$_T$ & Fun & neutral & GL & GL$\times$GL & M$\infty$ & -- & \\
$\bullet$\,Lily & \citealp{zhong2024low} & 2024 & F$_T$ & Fun & neutral & GL & -- & M$\infty$ & -- & \checkmark\\
LiME & \citealp{kowsher2026lime} & 2026 & F$_T$ & Fun & neutral & GL$\times$T & $\Pi$ & M$\infty$ & -- & \\
LoRA-Flow & \citealp{wang2024lora} & 2024 & F$_T$ & Fun & unpt. & none & -- & M$\infty$ & -- & \\
LoRA-Mixer & \citealp{li2025lora} & 2025 & F$_T$ & Fun & unpt. & GL & GL$\times$GL & M$\infty$ & -- & \\
$\bullet$\,LoRAMoE & \citealp{dou2023loramoe} & 2023 & F$_T$ & Fun & neutral & GL & -- & M$\infty$ & -- & \\
LoraRetriever & \citealp{zhao2024loraretriever} & 2024 & F$_T$ & Fun & unpt. & -- & -- & M$\infty$ & -- & \\
LPT & \citealp{liu2022late} & 2022 & F$_T$ & Fun & unpt. & torus & GL$\times$GL & M$\infty$ & -- & \\
LRAgent & \citealp{jeon2026lragent} & 2026 & F$_T$ & Cone & apex & GL & GL$\times$GL & M1 & -- & \\
MaLoRA / MaRA & \citealp{dokme2026selective} & 2026 & F$_T$ & Fun & neutral & GL & GL$\times$GL & M$\infty$ & -- & \\
$\bullet$\,MHR & \citealp{caccia2022multi} & 2022 & F$_T$ & Cone & apex & GL & $\Pi$ & M1 & span & \checkmark\\
MixLoRA & \citealp{li2024mixlora} & 2024 & F$_T$ & Fun & neutral & GL & -- & M$\infty$ & -- & \\
MOELoRA & \citealp{liu2023moe} & 2023 & F$_T$ & Cone & apex & GL & GL$\times$GL & M1 & span & \\
$\bullet$\,MoLE & \citealp{wu2024mixture} & 2024 & F$_T$ & Fun & unpt. & scalar & -- & M$\infty$ & -- & \\
MoV / MoLoRA & \citealp{zadouri2023pushing} & 2023 & F$_T$ & Fun & neutral & none & -- & M$\infty$ & -- & \\
NaRA & \citealp{wang2026nara} & 2026 & F$_T$ & Cone & apex & GL & GL$\times$GL & M1 & -- & \\
PERK & \citealp{chen2025perk} & 2025 & F$_T$ & Cone & defect & GL & GL$\times$GL & M1 & -- & \\
$\bullet$\,PHATGOOSE & \citealp{muqeeth2024learning} & 2024 & F$_T$ & Fun & unpt. & nonlin. & -- & M$\infty$ & -- & \\
$\bullet$\,Polytropon (Poly) & \citealp{ponti2022combining} & 2022 & F$_T$ & Cone & apex & GL & GL$\times$GL & M1 & span & \checkmark\\
PreFT & \citealp{lanpouthakoun2026preft} & 2026 & F$_T$ & Fun & apex & GL & GL$\times$GL & M$\infty$ & -- & \\
Prompt2Effect & \citealp{yang2026prompt2effect} & 2026 & F$_T$ & Cone & defect & GL & GL$\times$GL & M1 & -- & \\
$\bullet$\,Punica & \citealp{chen2023punica} & 2023 & F$_T$ & Cone & -- & GL & GL$\times$GL & M1 & -- & \\
RotMoLE & \citealp{sun2026rotmole} & 2026 & F$_T$ & Fun & neutral & GL & GL$\times$GL & M$\infty$ & -- & \\
$\bullet$\,S-LoRA & \citealp{sheng2023s} & 2023 & F$_T$ & Cone & -- & GL & GL$\times$GL & M1 & -- & \\
SHINE & \citealp{liu2026shine} & 2026 & F$_T$ & Cone & defect & GL & GL$\times$GL & M1 & -- & \\
SmartFed (MoRE + EEQA) & \citealp{wu2025elastic} & 2025 & F$_T$ & Fun & unpt. & none & GL$\times$GL & M$\infty$ & -- & \\
SMEAR & \citealp{muqeeth2023soft} & 2023 & F$_T$ & Fun & defect & torus & -- & M$\infty$ & -- & \\
$\bullet$\,Text-to-LoRA (T2L) & \citealp{charakorn2025text} & 2025 & F$_T$ & Cone & apex & GL & GL$\times$GL & M1 & -- & \\
VAPT & \citealp{le2025revisit} & 2025 & F$_T$ & Fun & unpt. & torus & GL$\times$GL & M$\infty$ & -- & \\
$\bullet$\,X-LoRA & \citealp{buehler2024x} & 2024 & F$_T$ & Fun & unpt. & none & -- & M$\infty$ & -- & \checkmark\\
\end{longtable}\endgroup


% generated by tools/make_tables.py
\begingroup\scriptsize\setlength{\tabcolsep}{3pt}\renewcommand{\arraystretch}{1.12}
\begin{longtable}{@{}>{\raggedright\arraybackslash}p{2.4cm}>{\raggedright\arraybackslash}p{2.4cm}c>{\raggedright\arraybackslash}p{8.6cm}@{}}
\caption{Every recorded simulation arrow $M\to N$ (rules A2--A5 of \S\ref{sec:atlas}), with its explicit map. Notes on the instances (rank, pointing) are in the ancillary file \texttt{arrows.csv}.}\label{tab:arrows}\\
\toprule\headrow\textbf{Source $M$} & \textbf{Target $N$} & \textbf{Rule} & \textbf{Pointed map $h$ with $\rho_N\circ h=\rho_M$}\\\midrule
\endfirsthead
\multicolumn{4}{@{}l}{\footnotesize\itshape Table~\ref{tab:arrows}, continued}\\
\toprule\headrow\textbf{Source $M$} & \textbf{Target $N$} & \textbf{Rule} & \textbf{Pointed map $h$ with $\rho_N\circ h=\rho_M$}\\\midrule
\endhead
\midrule
\multicolumn{4}{@{}r}{\footnotesize\itshape continued on the next page}\\
\endfoot
\bottomrule\endlastfoot
(IA)\textsuperscript{3} & DoRA & A2 & $\ell\mapsto\big(B=0,\ A=A_0,\ \mu=\ell\odot(\|w_{0,i}\|)_i\big)$\\
(IA)\textsuperscript{3} & HiRA & A3 & $\ell\mapsto(\ell-\mathbf 1,\ \mathbf 1^\top)\ \ (k,v);\qquad \ell\mapsto(\mathbf 1,\ (\ell-\mathbf 1)^\top)\ \ (W_{\rm down})$\\
ABBA & LoHa & A3 & $h(B_1,A_1,B_2,A_2)=(B_1,A_1,cB_2,A_2),\ \ c=\frac{\alpha/\sqrt{r_1r_2}}{\gamma_{\rm LoHa}}$\\
ABBA & LoRA & A3 & $h=\big(\tfrac{\alpha}{\sqrt{r_1r_2}}\,B_1\bullet B_2,\;(A_1^\top\bullet A_2^\top)^\top\big)$\\
AdaLomo & LOMO & A3 & $h=\mathrm{id}_\Theta$\\
AdaLoRA & LoRA & A3 & $h(P,\Lambda,Q)=(P\Lambda,\ Q)$\\
AdaLoRA & PoLAR & A3 & $h(P,\Lambda,Q)=\big(P(P^\top P)^{-1/2},\ (P^\top P)^{1/2}\Lambda(QQ^\top)^{1/2},\ Q^\top(QQ^\top)^{-1/2}\big)$\\
Adam-mini & LOMO & A3 & $h=\mathrm{id}_\Theta$\\
AdaMerging & LoRA Soups (CAT) & A3 & $h=\mathrm{id}:\ \lambda_i^{l}\mapsto\alpha_i^{l}$\\
AdaMix & Adapter (Houlsby) & A3 & $h\big((W^{down}_k)_{k\le M},\,W^{up}\big)=\big(\tfrac1M\textstyle\sum_kW^{down}_k,\ W^{up}\big)$\\
AdaMSS & LoRA & A3 & $h((C_k,B_k)_k)=\big([E_{\mathcal S_1}C_1\cdots E_{\mathcal S_K}C_K],\ [B_1\tilde A_1;\dots;B_K\tilde A_K]\big)$\\
AdaPaD & LoRA & A3 & $h=\big([\,b_{i,1}\ \cdots\ b_{i,r_i}\,],\ [\,a_{i,1}^\top;\dots;a_{i,r_i}^\top\,]\big)$\\
AdapterDrop & Adapter (Houlsby) & A2 & $h(q_{>n})=\big((U_l{=}0,\,D_l{=}D_l^0)_{l\le n},\ q_{>n}\big)$\\
AdapterDrop & Pfeiffer Adapter & A2 & $h(q_{>n})=\big((U_l{=}0,\,D_l{=}D_l^0)_{l\le n},\ q_{>n}\big)$\\
AdaptFormer & Parallel Adapter & A3 & $h(W_{down},b_{down},W_{up},b_{up})=\big(W_{down},b_{down},\tfrac{0.1}{s}W_{up},\tfrac{0.1}{s}b_{up}\big)$\\
AGZO & LoRA-FA & A3 & $\lambda\mapsto\lambda R_t$\\
ALoRA & LoRA & A3 & $h(B,A)=(B\,\mathrm{diag}(\alpha),\ A)$\\
AltLoRA & LoRA & A3 & $h=\mathrm{id}_Q$\\
AltLoRA & Riemannian LoRA & A3 & $h=\mathrm{id}_Q$\\
AP (Activation Prompt) & BitFit & A3 & $\delta=\mathbf 1\otimes v\ \mapsto\ \Delta b_{\rm down}^{(l-1)}=v$\\
AR1-ZO & LoRA & A3 & $h(B,A)=\big(\tfrac{\alpha}{s}B,\ A\big)$\\
Arrow & LoRA & A3 & $h_x(\ast)=\big(\textstyle\sum_ip_i(x)B_i,\ \sum_ip_i(x)A_i\big)$\\
Arrow & LoRA Soups (CAT) & A3 & $h_x(\ast)=\big(\alpha_i=p_i(x)\,s_i\big)_i$\\
Astra & LoRA & A3 & $h(B,A)=([B,\,-B_0],\,[A;\,A_0])$\\
AsyLoRA & LoRA & A2 & $h(B)=(B,\ Q)$\\
AsyLoRA & LoRA-FA & A3 & $h(B)=B\,QA_0^{+}$\\
AutoLoRA & LoRA & A3 & $h(b,a,\beta)=\big([\mathrm{softmax}(\beta)_jb_j]_j,\ [a_j^\top]_j\big)$\\
AutoPrompt & Prompt Tuning & A3 & $h(w_1,\dots,w_\ell)=(E_{w_1},\dots,E_{w_\ell})$\\
Bayesian-LoRA & LoRA & A3 & $h=\big(\bar B,\ \bar A\big),\ \ \bar B=T_r^B\,\mathbb E_q[U^B]\,T_c^B,\ \ \bar A=T_r^A\,\mathbb E_q[U^A]\,T_c^A$\\
BBT & Prompt Tuning & A3 & $h(z)=Az+p_0$\\
BD-LoRA & GraLoRA & A3 & $h(\{B_i\},A)=\big(B_{ij}=B_i,\;A_{ij}^\top=A^{(i)}_{:,J_j}\big)_{i,j\le N}$\\
BD-LoRA & LoRA & A3 & $h(B_1,\dots,B_N,A)=\big(\mathrm{diag}(B_1,\dots,B_N),\;A\big)$\\
BEFT & BitFit & A2 & $h(\Delta b_v)=\big(\Delta b_v,\,0_{\text{other biases}}\big)$\\
BiDoRA & DoRA & A3 & $\mathrm{id}$\\
BitFit & SSF & A2 & $h(\Delta b)=\big(\gamma=\mathbf 1,\ \beta=\Delta b\big)$\\
BLoB & LoRA & A3 & $h(M,G,B)=(B,\ M)$\\
BOFT & BOFT & A2 & $(Q_1,\dots,Q_{m'},s)\mapsto(Q_1,\dots,Q_{m'},0,s)$\\
Cartridges & Prefix-Tuning & A2 & $h(Z)=\big([z^{(l)}_{\mathrm{sink}};Z_k^{(l)}],\,[z^{(l)}_{v,\mathrm{sink}};Z_v^{(l)}]\big)_{l\le L}$\\
Child-Tuning & Diff Pruning & A2 & $x\mapsto(\alpha_0,\ \iota_M x)$\\
Child-Tuning & FISH Mask & A3 & $\mathrm{id}_{\mathbb R^M}$\\
CoA-LoRA & LQ-LoRA & A3 & $h_{\mathbf C}(\theta,L_1,L_2)=\big(L_1\,(I_r+U_\theta(Q_{\mathbf C})),\ L_2\big)$\\
COLA & LoRA & A3 & $h((B_i,A_i)_i)=([B_1\cdots B_M],\,[A_1;\dots;A_M])$\\
COLA & ReLoRA & A3 & $h=\mathrm{id}\ \text{(stagewise)}$\\
Compacter & Adapter (Houlsby) & A3 & $h=\big(W_{down}=\textstyle\sum_iA_i\otimes s_i^Dt_i^{D\top},\;W_{up}=\sum_iA_i\otimes s_i^Ut_i^{U\top}\big)$\\
Compacter & PHM-Adapter & A3 & $h\big(\{A_i\},\{s_i^{\bullet,\ell},t_i^{\bullet,\ell}\}\big)=\big(A_i^{\bullet,\ell}=A_i,\;B_i^{\bullet,\ell}=s_i^{\bullet,\ell}t_i^{\bullet,\ell\top}\big)_{\ell,\ \bullet\in\{D,U\}}$\\
Compress-then-Merge (CtM) & LoRA & A3 & $\lambda\mapsto\big(U\textstyle\sum_t\lambda_tO^{(t)},\ V^\top\big)$\\
Concept Sliders & LoRA & A3 & $h=\mathrm{id}$\\
Concept Sliders & Task Arithmetic & A3 & $h(\alpha_1,\dots,\alpha_K)=(\lambda_k=\alpha_k),\quad \tau_k=B_kA_k$\\
CoOp & CoCoOp & A2 & $h(V)=\big(V,\ \theta\ \text{with}\ W_2=0,\ b_2=0\big)$\\
CoOp & Prompt Tuning & A3 & $h=\mathrm{id}_{\mathbb R^{\ell\times d}}$\\
CorDA & LoRA & A3 & $h(B,A)=([B,\,-B_0],\,[A;\,A_0])$\\
Core Space & Task Arithmetic & A3 & $h=\mathrm{id}_\lambda$\\
CPT & Prompt Tuning & A2 & $h(E_{\mathrm{ctx}})=\big(E_{\mathrm{ctx}},\,E^{0}_{\mathrm{label}}\big)\quad(\|E_{\mathrm{ctx}}-E^0_{\mathrm{ctx}}\|\le\varepsilon)$\\
DC-Merge & LoRA & A3 & $h(\alpha)=\big(\alpha\,\widetilde U(\widetilde M\odot\mathcal M),\ \widetilde V^\top\big)$\\
DEFT & LoRA & A3 & $(P,R)\mapsto\big(P,\ R-\mathrm{relu}(P)^\top W_0\big)\quad(\text{QR: }(Q(P),\ R-Q(P)^\top W_0))$\\
DeLoRA & DoRA & A3 & $(B,A,\lambda)\mapsto\big(B',A',\ \mu=(\|(W_0+B'A')_{i,:}\|)_i\big),\ (B',A')\ \text{as in the LoRA arrow}$\\
DeLoRA & LoRA & A3 & $(B,A,\lambda)\mapsto\Big(\tfrac{\lambda}{r}\,B\,\mathrm{diag}\big(1/\max(\|b_i\|,\varepsilon)\big),\ \mathrm{diag}\big(1/\|a_i\|\big)\,A\,D_w\Big)$\\
Delta-LoRA & LoRA & A3 & $h(B,A)=\big([(1+\lambda)B,\ -\lambda B_K],\ [A;\ A_K]\big)$\\
DiaBlo & Diff Pruning & A2 & $D\mapsto\big(\alpha_0,\ \iota_{\rm blk}D\big)$\\
DiaBlo & FISH Mask & A3 & $h=\mathrm{id},\qquad M=\mathrm{supp}\,\mathrm{blkdiag}(\mathbf 1_{d_{out}/N\times d_{in}/N},\dots)$\\
DnD & LoRA & A3 & $h_c(\phi)=\mathrm{detok}\big(G_\phi(c)\big)$\\
Doc-to-LoRA (D2L) & LoRA & A3 & $h_c(\phi)=\big(\alpha_l[B^{(1)}_l\ \cdots\ B^{(K)}_l],\ [A^{(1)}_l;\dots;A^{(K)}_l]\big)$\\
DoRA & BiDoRA & A3 & $\mathrm{id}$\\
DoRA & GLoRA & A3 & $(B,A,\mu)\mapsto\big(B_A=c-\mathbf 1,\ A_A=\mathbf 1^\top;\ B_B=\mathrm{diag}(c)\,sB,\ A_B=A\big),\quad c=\mu\oslash\big(\|(W_0+sBA)_{i,:}\|\big)_i$\\
DoRA (dynamic rank) & AdaLoRA & A3 & $h=\mathrm{id}$\\
DoRA (dynamic rank) & LoRA & A3 & $h(B,c,A)=(B\,\mathrm{diag}(c),\ A)$\\
DoRA (dynamic rank) & SoRA & A3 & $h(B,c,A)=\big(B\,\mathrm{diag}(c)\,\mathrm{diag}(g_0)^{-1},\ g_0,\ A\big)$\\
DSEE & LoRA & A3 & $(U,V,S)\mapsto\big([\,U\mid E_I\,\mathrm{diag}(S_\Omega)\,],\ [\,V;\ E_J^\top\,]\big),\quad E_I=[e_{i_1}\cdots e_{i_c}],\ E_J=[e_{j_1}\cdots e_{j_c}]$\\
DyLoRA & DyLoRA & A2 & $h(B_b,A_b)=\big([\tfrac{r_{\max}}{b}B_b,\ 0],\ [A_b;\ A_{0,(b+1):}]\big)$\\
DyLoRA & LoRA & A3 & $h=\mathrm{id}$\\
EasyControl & Prefix-Tuning & A3 & $h_c(B,A)=\big(K_l(c;B,A),\ V_l(c;B,A)\big)_{l\le L}$\\
EFT & Proxy-Tuning & A3 & $h=\mathrm{id}$\\
ElaLoRA & AdaLoRA & A3 & $h=\mathrm{id}$\\
ElaLoRA & LoRA & A3 & $h(P,\Lambda,Q)=(P\Lambda,\ Q)$\\
ETHER+ & LoRA & A3 & $q\mapsto\Big(\big[\,[\hat{\tilde v}_j,-\hat{\tilde u}_j]_j,\ W_0[\hat v_j,-\hat u_j]_j\,\big],\ \big[\,[\hat{\tilde v}_j,\hat{\tilde u}_j]_j^\top W_0H^+;\ [\hat v_j,\hat u_j]_j^\top\,\big]\Big)$\\
EVA & LoRA & A3 & $h=\mathrm{id}:(B,A)\mapsto(B,A)$\\
FacT & FacT & A3 & $h(U,P,C,V)=\big(U,\ (\textstyle\sum_tP_{i,t}C_t)_i,\ V\big)$\\
FacT & LoRA & A3 & $h(U,\Sigma,V)=\big(sU\Sigma_i,\ V^\top\big)_i$\\
FedRot-LoRA & LoRA & A3 & $h=\mathrm{id}$\\
FFA-LoRA & LoRA & A2 & $h(B)=(B,\ A_0)$\\
FFA-LoRA & LoRA-FA & A3 & $h=\mathrm{id}$\\
FISH Mask & Diff Pruning & A2 & $x\mapsto(\alpha_0,\ \iota_M x)$\\
FISH Mask & RoSA & A5 & $x\mapsto\big((0,A_0),\ x\big)$\\
Flat-LoRA & LoRA & A3 & $h=\mathrm{id}_Q$\\
FlexLoRA & AdaLoRA & A3 & $h=\mathrm{id}$\\
FlexLoRA & ElaLoRA & A3 & $h=\mathrm{id}$\\
FlexLoRA & LoRA & A3 & $h(P,\Lambda,Q)=(P\Lambda,\ Q)$\\
FLoRA & HiRA & A3 & $h(B,A)=\big([\tfrac1rB\;\,{-\mathbf 1_m}],\,[A;\mathbf 1_n^\top]\big)$\\
Flora & LoRA-FA & A3 & $C\mapsto s^{-1}C\quad(A_0:=A_t)$\\
Flora & ReLoRA & A2 & $C\mapsto(s^{-1}C,\ A_t)$\\
FLoRA (federated) & ReLoRA & A3 & $h\big((B_k,A_k)_k\big)=\big([p_1B_1\cdots p_KB_K],\ [A_1;\dots;A_K]\big)\ \ \text{per round}$\\
FLoRA (Tucker core) & LoRA & A3 & $h(A,G,B)=\big(sAG,\;B^\top\big)$\\
FLoRG & LoRA & A3 & $A\mapsto\big(s\,LA^\top,\ AR\big)$\\
FLoRG & SingLoRA & A3 & $A\mapsto\sqrt{sr/\alpha}\,A^\top$\\
FourierFT & FourierFT & A2 & $h(c_\Omega)=(c_\Omega,\ 0_{\Omega'\setminus\Omega})$\\
FourierFT & LoRA-FA & A3 & $h(c)=\tfrac{\alpha}{mn}\,B_F\,\mathrm{diag}(c,-c),\ \ B_F=[\cos\theta_k\ \sin\theta_k]_k,\ A_F=[\cos\varphi_k;\ \sin\varphi_k]_k$\\
FourierMoE & LoRA-FA & A3 & $h_x\big((c_i)_i,\Phi\big)=\eta\textstyle\sum_{i\in\mathcal S_\Phi(x)}g_{\Phi,i}(x)\,B(c_i),\ \ B(c)=\tfrac{2}{d_{out}d_{in}}\big[\,a_k\cos\theta_k-b_k\sin\theta_k,\ -(a_k\sin\theta_k+b_k\cos\theta_k)\,\big]_k$\\
FrameFT & LoRA-FA & A3 & $h(C_\Omega)=\tfrac{\alpha}{\sqrt{d_{out}d_{in}}}\,P_{out}\,C_{:,J},\qquad A_0=(P_{in})_{:,J}^\top$\\
FRoD & FRoD & A2 & $h(\Sigma)=(\Sigma,\ S=0)$\\
GaLore & LoRA & A2 & $C\mapsto(P_k,\ s^{-1}C)$\\
GaLore & LoRA-FA & A3 & $C\mapsto s^{-1}C\quad(A_0:=P_k^\top)$\\
GaLore 2 & GaLore & A3 & $\mathrm{id}$\\
GenerativeAdapter & LoRA & A3 & $h_c(A_1,A_2,B_1,B_2)=\big(A_1U_r(c),\ V_r(c)^\top B_2\big),\ \ U_r\Sigma_rV_r^\top=\mathrm{SVD}_r(A_2H_c^\top H_cB_1)$\\
GeoRA & LoRA & A3 & $(B,A)\mapsto\big([B,\ -B_0],\ [A;\ A_0]\big)$\\
GiVA & LoRA & A3 & $(\gamma,\lambda)\mapsto(\Gamma B\Lambda,\ V_r^\top)$\\
GiVA & LoRA-FA & A3 & $(\gamma,\lambda)\mapsto\Gamma B\Lambda$\\
GOAT & LoRA & A3 & $h_x(B,A,W_z)=\big([s\,w^1(x)B^1\ \cdots\ s\,w^E(x)B^E\ \ -\tfrac sEB_0^1\ \cdots\ -\tfrac sEB_0^E],\ [A^1;\dots;A^E;A_0^1;\dots;A_0^E]\big)$\\
GOFT & BOFT & A3 & $\theta\mapsto\big(t=\tan(\theta/2)\ \text{on the kept }2\times2\text{ blocks},\ 0\ \text{on the others},\ s=\mathbf 1\big)$\\
GOFT & qGOFT & A3 & $\theta\mapsto\Big(\tilde G_t=\begin{pmatrix}\cos\theta_t&-\sin\theta_t\\ \sin\theta_t&\cos\theta_t\end{pmatrix}\Big)_t$\\
GoRA & LoRA & A3 & $h(B,A)=(cB,A),\ c=\tfrac{\alpha/\sqrt{r_\ell}}{s}$\\
GPS & Diff Pruning & A2 & $x\mapsto(\alpha_0,\ \iota_M x)$\\
GradFix & FISH Mask & A3 & $h(\alpha)=\alpha\,(m\odot\tau_A)$\\
GraLoRA & LoRA & A3 & $h(\{U_{ij},V_{ij}\})=\big(\mathrm{blockdiag}(L_1,\dots,L_k),\;S\big),\ L_i=[U_{i1}\cdots U_{ik}],\ S=\text{stack of the }V_{ij}^\top\text{ on column block }j$\\
HetLoRA & DyLoRA & A3 & $h=\mathrm{id}$\\
HetLoRA & LoRA & A3 & $h=\mathrm{id}$\\
HiFT & Surgical FT & A3 & $\mathrm{id}_{\Theta_{G_j}}$\\
HiRA & GLoRA & A2 & $h(B,A)=\big(\mathcal A\text{-path}=BA,\ \text{all other GLoRA paths frozen at }0\big)$\\
HiRA & LoHa & A2 & $h(B,A)=(U_0,V_0,B,A),\ \ U_0V_0=W_0/\gamma\ \text{a fixed rank-}\mathrm{rk}W_0\text{ factorisation}$\\
HRA & LoRA & A3 & $u\mapsto\big(W_0B(u),\ A(u)\big),\quad B(u)=\big[-2H_{u_1}\cdots H_{u_{i-1}}\hat u_i\big]_{i=1}^r,\ A(u)=\big[\hat u_1^\top;\dots;\hat u_r^\top\big]$\\
HydraLoRA & LoRA & A3 & $h_x(A,B_1,\dots,B_N,W_g)=\big(\textstyle\sum_i\omega_i(x)B_i,\ A\big)$\\
HydraLoRA & LoRAMoE & A3 & $h(A,(B_i)_i,W_g)=\big((\tfrac r\alpha B_i)_i,\ (A,\dots,A),\ W_g\big)$\\
HyperDreamBooth & LoRA & A3 & $h(B_{tr},A_{tr})=\big(P_{out}B_{tr},\ A_{tr}P_{in}\big)$\\
HyperDreamBooth & NOLA & A3 & $h(B_{tr},A_{tr})=\big(\alpha=\mathrm{vec}\,A_{tr},\ \beta=\mathrm{vec}\,B_{tr}\big),\ \text{NOLA bases }\{p_ke_j^\top\},\{e_jq_l^\top\}$\\
IC-LoRA & LoRA & A3 & $h=\mathrm{id}$\\
ICEdit & LoRA & A3 & $h_x(B,A,g)=\big(\tfrac\alpha r[G_1(x)B_1\ \cdots\ G_4(x)B_4],\ [A_1;\dots;A_4]\big)$\\
In-Place TTT & LoRA & A3 & $h_{c,i}(W_{target},k)=\big(\eta\,[\hat V_{[1]}^\top\ \cdots\ \hat V_{[i-1]}^\top],\ [Z_{[1]};\dots;Z_{[i-1]}]\big)$\\
IncreLoRA & AdaLoRA & A3 & $h=\mathrm{id}:(B,\Lambda,A)\mapsto(P,\Lambda,Q)$\\
IncreLoRA & IncreLoRA & A5 & $h(q)=\big(q,\ (b_{\rm new},0,a_{\rm new})\big)$\\
IncreLoRA & LoRA & A3 & $h(B,\Lambda,A)=(B\Lambda,\ A)$\\
InfLoRA & LoRA & A2 & $h(B)=(B,\ A_t)$\\
InfLoRA & LoRA-FA & A3 & $h=\mathrm{id}$\\
Initial State Tuning & Prefix-Tuning & A3 & $h(h')=\big(\mathcal C_{\ell,l}^{+}\,h'_l\big)_l,\quad \mathcal C_{\ell,l}=[\bar B_l,\bar A_l\bar B_l,\dots,\bar A_l^{\ell-1}\bar B_l]$\\
Intruder-dimension reduction & SVDiff & A3 & $h(\lambda)=\delta,\quad \delta_j=(\lambda-1)\sigma_j\ (j\in\mathcal I),\ \ \delta_j=0\ \text{otherwise}$\\
IR-QLoRA & QLoRA & A3 & $(B,A,\beta_1,\beta_2)\mapsto(\tfrac{\alpha}{s}(B+\beta_2E),\ A+\beta_1P)$\\
K-LoRA & LoRA Soups (CAT) & A3 & $h_t(\ast)=\big(\alpha_c^{l},\alpha_s^{l}\big)=\begin{cases}(1,0)&\text{layer }l\text{ selects content at }t\\ (0,1)&\text{otherwise}\end{cases}$\\
K-LoRA & ZipLoRA & A3 & $h_t(\ast)=\big(m_c^{l},m_s^{l}\big)\in\{(\mathbf 1,0),(0,\mathbf 1)\}$\\
KAdaptation & LoRA & A3 & $h=\big([A_1\otimes u_1,\dots,A_n\otimes u_n],\;[I_n\otimes v_1^\top;\dots;I_n\otimes v_n^\top]\big)$\\
KaSA & LoRA & A3 & $h(B,\Delta\sigma,A)=\big(\tfrac{\alpha}{r}B\,\mathrm{diag}(\Delta\sigma),\ A\big)$\\
KeepLoRA & LoRA & A3 & $B\mapsto\big(B-B_t^{(0)},\ A_t\big)$\\
KeepLoRA & LoRA-FA & A3 & $B\mapsto B-B_t^{(0)}$\\
KnOTS & Task Arithmetic & A3 & $h=\mathrm{id}_\lambda$\\
KRAdapter & LoHa & A2 & $h(U,V)=(B_1=S_1,\;A_1=U,\;B_2=S_2,\;A_2=V),\ \ S_1,S_2\ \text{the }\{0,1\}\text{ row selectors of the Kronecker index}$\\
KronA & LoKr & A3 & $h(A_k,B_k)=\big(C=A_k,\;\text{right factor}=B_k\big)\ \ \text{(LoKr with neither factor decomposed)}$\\
LaMDA & LoRA & A3 & $h(U,S)=\big(\alpha US-U_r,\;\Sigma_rV_r^\top\big)$\\
LaMDA & LoRA-FA & A3 & $h(U,S)=\alpha US-U_r\ \text{ into }\ \mathrm{FA}(\Sigma_rV_r^\top)$\\
LaMDA & PiSSA & A3 & $h(U,S)=\big(\alpha US\,\Sigma_r^{1/2},\;\Sigma_r^{1/2}V_r^\top\big)$\\
Laplace-LoRA & LoRA & A3 & $h=\mathrm{id}\ \ \text{(the MAP object)}$\\
LCM-LoRA & LoRA & A3 & $h(\lambda_1,\lambda_2)=\big([\lambda_1B'\ \ \lambda_2B_{\rm LCM}],\ [A';A_{\rm LCM}]\big)$\\
LCM-LoRA & Task Arithmetic & A3 & $h(\lambda_1,\lambda_2)=(\lambda_1,\lambda_2),\quad \tau_1=B'A',\ \tau_2=B_{\rm LCM}A_{\rm LCM}$\\
LD-MoLE & LoRA & A3 & $h_x\big((B_i,A_i)_i,W_{gate},f\big)=\big(\tfrac{\alpha}{rs}\,[\,p_{x,1}B_1\ \cdots\ p_{x,E}B_E\,],\ [A_1;\dots;A_E]\big)$\\
LieRA & FLoRA & A3 & $h(B,A)=\big([\mathbf 1_m\;B],\,[\mathbf 1_n^\top;A]\big)$\\
LieRA & HiRA & A3 & $h(B,A)=(B,A)$\\
LieRA & LoHa & A3 & $h(B,A)=(B,A,U_0,V_0^\top),\quad U_0V_0^\top=W_0$\\
LIFT & Diff Pruning & A2 & $x\mapsto(\alpha_0,\ \iota_{M_0}x)$\\
LiME & LoRA & A3 & $h_x(B,A,(p_i)_i,p_s,\gamma)=\big(\mathrm{diag}(\mathcal P(x)+\gamma p_s)\,B,\ A\big)$\\
Linear probing & BitFit & A2 & $h(V)=(\Delta b=0,\ V)$\\
Linear probing & LP-FT & A5 & $V\mapsto(V,\ q_0^{(2)})$\\
Linear probing & Surgical FT & A2 & $h(V)=V\ \ \text{(block = last layer)}$\\
LLaMA-Adapter & LLaMA-Adapter V2 & A2 & $h(P,g)=\big(P,\,g,\ s=\mathbf 1,\ b=0,\ \gamma=\gamma_0,\ \text{visual branch at its zero gate}\big)$\\
LN Tuning & SSF & A3 & $(\gamma,\beta)\mapsto\big(\gamma'=\gamma\oslash\gamma_0,\ \beta'=\beta-(\gamma\oslash\gamma_0)\odot\beta_0\big)\ \text{on SSF's post-norm wrappers; other wrappers at }(\mathbf 1,0)$\\
LoCA & LoRA-FA & A3 & $h(a)=\big[a_1C_{l^1_1,:}^\top\ \cdots\ a_{\mathcal B}C_{l^1_{\mathcal B},:}^\top\big]\ \text{ into }\ \mathrm{FA}(A_0),\ A_0=\big[D_{l^2_1,:};\dots;D_{l^2_{\mathcal B},:}\big]$\\
LoCO & LoRA & A3 & $(U,V)\mapsto\big(2W_0X(I_{2r}-Y^\top X)^{-1},\ Y^\top\big),\ \ X=[U\,|\,-V],\ Y=[V\,|\,U]$\\
LoCO & OFT & A3 & $(U,V)\mapsto Q=-(UV^\top-VU^\top)$\\
LoDA & LoRA & A3 & $(B_G,B_I)\mapsto\big([\,w_GB_G,\ B_I\,],\ [A_G;A_I]\big)$\\
LoDA & LoRA-FA & A3 & $(B_G,B_I)\mapsto[\,w_GB_G,\ B_I\,]$\\
LoFiT & BitFit & A3 & $h(v)=\big(\Delta b_O^{(l)}=\textstyle\sum_{i\in T_l}W^{O}_{(i)}v_l^i,\ \text{other biases }0\big)$\\
LoFT & LoRA & A3 & $h=\mathrm{id}_Q$\\
LOFT & LoRA & A3 & $E\mapsto\big(B=W_0P_r^{\top}(C(E)-I_r),\ A=P_r\big)$\\
LOFT & LoRA-XS & A3 & $E\mapsto R=\Sigma_r\big(C(E)-I\big)\Sigma_r^{-1}$\\
LOFT & OFT & A3 & $E\mapsto Q=-P_r^{\top}EP_r\in\mathfrak{so}(n)$\\
LOFT & PSOFT & A3 & $E\mapsto\big(R=C(E)^{\top},\ \alpha=\beta=\mathbf 1\big)$\\
LoftQ & LoRA & A3 & $h(B,A)=([B,\,-B_T],\,[A;\,A_T])$\\
LoHa & LoRA & A3 & $h(B_1,A_1,B_2,A_2)=\big(\gamma\,B_1\bullet B_2,\;(A_1^\top\bullet A_2^\top)^\top\big)$\\
LoKr & KronA & A3 & $h(C,B,A)=(C,\;BA)$\\
LoKr & LoRA & A3 & $h(C,B,A)=\big(C\otimes B,\;I_{u_q}\otimes A\big)$\\
LoQT & GaLore & A3 & $B\mapsto sB$\\
LoQT & LoRA & A2 & $B\mapsto(P_k,\ \tfrac{s}{s^\prime}B)$\\
LoRA & DoRA & A3 & $h(B,A)=\big((\lVert(W_0+sBA)_{i,:}\rVert)_i,\ B,\ A\big)$\\
LoRA & LoRA & A3 & $h(B,A)=([B\ 0],\,[A;\,a_0])$\\
LoRA & PEANuT & A3 & $h(B,A)=\big([M,\,-M],\ [B^\top;\,-B^\top]\big),\quad M=(W_0^{+})^\top A^\top$\\
LoRA & ReLoRA & A5 & $q\mapsto(q,\,q_0^{(2)},\dots,q_0^{(K)})$\\
LoRA Dropout & LoRA & A3 & $h=\mathrm{id}_Q$\\
LoRA Soups (CAT) & AdaMerging & A3 & $h=\mathrm{id}:\ \alpha_i^{l}\mapsto\lambda_i^{l}$\\
LoRA Soups (CAT) & LoRA & A3 & $h(\alpha)=\big([\alpha_1B_1\ \cdots\ \alpha_NB_N],\ [A_1;\dots;A_N]\big)$\\
LoRA Switch / LoRA Composite & LoRA Soups (CAT) & A3 & $h_t(\ast)=w_{i(t)}\,e_{i(t)}$\\
LoRA Without Regret & LoRA & A3 & $h=\mathrm{id}_Q$\\
LoRA+ & LoRA & A3 & $h=\mathrm{id}_Q$\\
LoRA-Dash & LoRA & A3 & $h(B,A,\Delta\sigma)=\big([B\ \ \bar U\mathrm{diag}(\Delta\sigma)],\ [A;\ \bar V^\top]\big)$\\
LoRA-drop & LoRA & A3 & $h\big((B_i,A_i)_{\rm keep},(B_s,A_s)\big)=\big((B_i,A_i)_{\rm keep},(B_s,A_s)_{j\in\rm share}\big)$\\
LoRA-FA & AsyLoRA & A3 & $h(B)=B\,A_0Q^\top$\\
LoRA-FA & Flora & A5 & $B\mapsto(B,\,0,\dots,0)$\\
LoRA-FA & LoRA & A2 & $h(B)=(B,\ A_0)$\\
LoRA-Flow & LoRA Soups (CAT) & A3 & $h_x(W_{\rm gate},b)=\big(\alpha_i^{l}=w_i^{l}(x)\big)_{i,l}$\\
LoRA-GA & LoRA & A3 & $h(B,A)=([B,\,-B_0],\,[A;\,A_0])$\\
LoRA-LEGO & LoRA Soups (CAT) & A3 & $h(\ast)=N^{-1/2}\mathbf 1$\\
LoRA-Mini & LoRA & A3 & $h(C,D)=(PC,\;DQ)$\\
LoRA-Mixer & LoRA & A3 & $h_x\big((B_e,A_e)_e,G\big)=\big(\,[\,\mathbb 1_{e\in\mathrm{TopK}(x)}\,p_e(x)\,B_e\,]_{e=1}^{E},\ [A_1;\dots;A_E]\,\big)$\\
LoRA-Null & LoRA & A3 & $h(B,A)=([B,\,-B_0],\,[A;\,A_0])$\\
LoRA-Pre & Adam-mini & A3 & $h=\mathrm{id}_\Theta$\\
LoRA-Pro & LoRA & A3 & $h=\mathrm{id}_Q$\\
LoRA-RITE & LoRA & A3 & $h=\mathrm{id}_Q$\\
LoRA-RLPO / LoRA-RLMO & LoRA & A3 & $h=\mathrm{id}:(B,A)\mapsto(B,A)$\\
LoRA-S (AdamS, LRACS) & LoRA & A3 & $h=\mathrm{id}_Q$\\
LoRA-SB & LoRA-FA & A4 & $h(R)=sU_rR$\\
LoRA-X & LoRA & A3 & $h(\Delta\Sigma)=(U_r,\ \Delta\Sigma V_r^\top)$\\
LoRA-X & LoRA-XS & A3 & $h(\Delta\Sigma)=\Delta\Sigma\,\Sigma_r^{-1}$\\
LoRA-X & SVDiff & A2 & $h(\delta\sigma)=(\delta\sigma_1,\dots,\delta\sigma_r,0,\dots,0)$\\
LoRA-X & SVFT & A2 & $h(\delta\sigma)=M,\ \ M_{ii}=\delta\sigma_i\ (i\le r),\ \ 0\ \text{otherwise}$\\
LoRA-XS & LoRA & A3 & $h(R)=\big(U_rR,\;\Sigma_rV_r^\top\big)$\\
LoRA-XS & LoRA-FA & A4 & $h(R)=U_rR\ \text{ into }\ \mathrm{FA}(A_0),\ A_0=\Sigma_rV_r^\top$\\
LoRA-XS & PiSSA & A3 & $h(R)=\big(U_r\Sigma_r^{1/2}+U_rR\,\Sigma_r^{1/2},\;\Sigma_r^{1/2}V_r^\top\big)$\\
LoRA-XS & SVFT & A3 & $h(R)=\begin{bmatrix}R\Sigma_r&0\\0&0\end{bmatrix}\ \ (\Omega=\text{top-left }r\times r\text{ block})$\\
LoraHub & LoRA & A3 & $h(w)=\big(\textstyle\sum_iw_iB_i,\ \sum_iw_iA_i\big)$\\
LoRAMoE & LoRA & A3 & $h_x(B,A,W_g)=\big(\tfrac\alpha r[\omega_1(x)B_1\ \cdots\ \omega_N(x)B_N],\ [A_1;\dots;A_N]\big)$\\
LoRAMoE & MoV / MoLoRA & A3 & $h(B_i,A_i,W_g)=(\tfrac\alpha rB_i,\,A_i,\,W_g^\top)$\\
LoRAPrune & DoRA & A3 & $h(B,A,m)=\big(\mu=m\odot\lVert W_0+sBA\rVert_{\rm row},\ B,\ A\big)$\\
LoraRetriever & LoRA Soups (CAT) & A3 & $h_x(\ast)=\big(\alpha_\phi=\mathbf 1[\phi\in\Phi_x]/k\big)_\phi$\\
LoraRetriever & LoraHub & A3 & $h_x(\ast)=w=\mathbf 1_{\Phi_x}/k$\\
LoReFT & LoRA & A3 & $h(R,W,b)=\big(R^\top,\ [W-R\mid b]\big)$\\
LoRETTA & LoRA & A3 & $h(\{G^i_{up}\},\{G^i_{down}\})=\big(\mathrm{TT}(\{G^i_{up}\}),\;\mathrm{TT}(\{G^i_{down}\})\big)$\\
LoRI & LoRA & A2 & $h(X)=(\iota_M X,\ A_0)$\\
LoRI & LoRA-FA & A2 & $h(X)=\iota_M X$\\
LoRTA & FacT & A3 & $h=\big(U=\mathcal B,\;V=A,\;\Sigma_{(l,i)}=\mathrm{Diag}(C_L[l,:]\odot C_M[i,:])\big)$\\
LoRTA & LoRA & A3 & $h=\big(\mathcal B\,\mathrm{Diag}(C_L[l,:]\odot C_M[i,:]),\;A^\top\big)_{(l,i)},\ \ \mathcal B=\big[B\,\mathrm{Diag}(C_H[1,:]);\dots;B\,\mathrm{Diag}(C_H[H,:])\big]$\\
LoRTA & LoTR & A3 & $h=\big(A_{\rm LoTR}=A,\;B_{\rm LoTR}=\mathcal B,\;G_{(l,i)}^\top=\mathrm{Diag}(C_L[l,:]\odot C_M[i,:])\big)$\\
LoRTA & MetaTT & A3 & $h=\big(G_1=A,\;G_2[l]=\mathrm{Diag}(C_L[l,:]),\;G_3[i]=\mathrm{Diag}(C_M[i,:]),\;G_H[k]=\mathrm{Diag}(C_H[k,:]),\;G_4=B^\top\big)$\\
LoTR & FacT & A3 & $h(A,B,\{G_s\})=\big(V=A,\;U=B,\;\Sigma_s=G_s^\top\big)$\\
LoTR & FLoRA (Tucker core) & A3 & $h(A,B,\{G_s\})=\big(A_F=B,\;G_F=G_s^\top,\;B_F=A\big)_s$\\
LoTR & LoRA & A3 & $h=\big(BG_s^\top,\;A^\top\big)_s$\\
LoTR & ShareLoRA & A3 & $h=\big(\{B_s=BG_s^\top\}_s,\;A^\top\big)$\\
LRAgent & LoRA & A3 & $h_i(A,B_1,\dots,B_N)=(B_i,\ A)$\\
LRAgent & S-LoRA & A3 & $(A,(B_i)_i)\mapsto\big((B_i,A)\big)_i$\\
LT-SFT & Diff Pruning & A2 & $x\mapsto(\alpha_0,\ \iota_M x)$\\
MaLoRA / MaRA & LoRA & A3 & $h_{x_{1:t}}(A,B,P,\mathrm{Mamba},w,b)=\big(\tfrac{\alpha}{rs}\,\lambda(x_{1:t})\,B,\ A\big)$\\
MatryoshkaLoRA & DyLoRA & A3 & $h=\mathrm{id}$\\
MatryoshkaLoRA & LoRA & A3 & $(B,A)\mapsto\big(s_kB_{:,1:k},\ A_{1:k,:}\big)$\\
MatryoshkaLoRA & SoRA & A2 & $(B,A)\mapsto(B,\ p,\ A)$\\
MDPD & LN Tuning & A3 & $h=\mathrm{id}$\\
MELoRA & GraLoRA & A2 & $(B_i,A_i)_i\mapsto\big(B_{ii}=B_i,\ A_{ii}=A_i^\top;\ (B_{ij},A_{ij})=(0,A_{ij,0})\ (i\neq j)\big)$\\
MELoRA & LoRA & A2 & $(B_i,A_i)_i\mapsto\big(\mathrm{blockdiag}(B_1,\dots,B_n),\ \mathrm{blockdiag}(A_1,\dots,A_n)\big)$\\
MetaTT & LoRA & A3 & $h=\big(G_4^\top G_3[m]^\top G_2[l]^\top,\;G_1^\top\big)_{(l,m)}$\\
MetaTT & LoTR & A3 & $h(G_1,\{G_2[l]\},\{G_3[m]\},G_4)=\big(A=G_1,\;B=G_4^\top,\;G_{(l,m)}=G_2[l]G_3[m]\big)\ \text{per type }m$\\
MHR & LoRA & A3 & $h_\tau(A,B,Z)=(B^\tau,\,A^\tau)$\\
MiCA & LoRA & A2 & $h(A)=\big(U_{-r},\ \tfrac{\alpha}{r}A\big)$\\
MiCA & MiLoRA & A3 & $h(A)=\big(U_m\Sigma_m^{1/2},\ \Sigma_m^{1/2}V_m^\top+\tfrac{\alpha}{r}\Sigma_m^{-1/2}A\big)$\\
MiLoRA & LoRA & A3 & $h(B,A)=([B,\,-B_0],\,[A;\,A_0])$\\
MiSS (formerly Bone) & LoRA & A2 & $h(D)=(D,\ \mathbf 1^\top_{d_{\rm in}/r}\otimes I_r)$\\
MiSS (formerly Bone) & LoRA-FA & A3 & $h(D)=D,\ \ A_0=\mathbf 1^\top_{d_{\rm in}/r}\otimes I_r$\\
Mix-of-Show & RegMean & A3 & $h(\ast)=(\alpha=1)$\\
Model Soups & Task Arithmetic & A3 & $h(\ast)=\mathbf 1_S/|S|$\\
ModuLoRA & QLoRA & A3 & $h=\mathrm{id}$\\
MOELoRA & LoRA & A3 & $h_\tau(B,A,e,W_T)=\big([\omega_{\tau1}B_1\ \cdots\ \omega_{\tau N}B_N],\ [A_1;\dots;A_N]\big)$\\
MonteCLoRA & LoRA & A3 & $h(B,\mu,\xi)=(B,\ \mu)$\\
MoRA & KronA & A2 & $M\mapsto(A_k,B_k)=(I_{d/\hat r},\ M/s)$\\
MoRe & LoRA & A2 & $(L,R)\mapsto(P_1LP_2,\ R)$\\
MoS & LoRA & A3 & $h(\mathrm{pool}_B,\mathrm{pool}_A)=\big(\mathrm{Route}^c(\mathrm{pool}_B,I^b_\ell),\;\mathrm{Route}^r(\mathrm{pool}_A,I^a_\ell)\big)_\ell$\\
MoSLoRA & FLoRA (Tucker core) & A3 & $h=\mathrm{id}$\\
MoSLoRA & LoRA & A3 & $h(B,M,A)=(BM,\ A)$\\
MoV / MoLoRA & (IA)\textsuperscript{3} & A3 & $h_x(\ell_1,\dots,\ell_N,W_g)=\textstyle\sum_is_i(x)\ell_i$\\
MoV / MoLoRA & LoRAMoE & A3 & $h(B_i,A_i,W_g)=(\tfrac r\alpha B_i,\,A_i,\,W_g^\top)\ \text{(MoLoRA variant)}$\\
MPT & Prompt Tuning & A3 & $h(P^*,u,v)=P^*\odot uv^\top$\\
NaRA & LoRA & A3 & $h_\lambda(B,A,\phi)=\big(B\,C_\phi(\lambda),\ A\big)$\\
NaRA & MoSLoRA & A3 & $h_\lambda(B,A,\phi)=\big(B,\ C_\phi(\lambda)/s,\ A\big)$\\
NOLA & LoRA & A3 & $h(\alpha,\beta)=\Big(\sum_j\beta_jB_j,\ \sum_i\alpha_iA_i\Big)$\\
O-LoRA & LoRA & A5 & $h=\big([B_1\ \cdots\ B_T],\ [A_1;\dots;A_T]\big)$\\
O-LoRA & ReLoRA & A3 & $h=\mathrm{id}\ \text{on each stage's }(B_t,A_t)$\\
OFT & BOFT & A2 & $Q\mapsto(Q,\ s=\mathbf 1)\ \text{in BOFT}(1,b)$\\
OFT & GOFT & A3 & $t\mapsto\big(\theta^{(1)}_k=2\arctan t_k,\ \theta^{(\ell)}=0\ (\ell\ge2)\big)$\\
OFTv2 & qGOFT & A3 & $t\mapsto\Big(\tilde G^{(1)}_k=(1-t_k^2)\begin{pmatrix}1-t_k^2&-2t_k\\2t_k&1-t_k^2\end{pmatrix},\ \tilde G^{(\ell)}=I\ (\ell\ge2)\Big)$\\
OLoRA & LoRA & A3 & $h(B,A)=([B,\,-B_0],\,[A;\,A_0])$\\
Online Subspace Descent & LoRA & A2 & $h(X)=(P_t,\,X)$\\
Orthogonal Adaptation & LoRA & A2 & $h(B)=(B,\,A_i)$\\
Orthogonal Adaptation & LoRA-FA & A3 & $h=\mathrm{id}$\\
OrthoMerge & OFT & A3 & $\ast\mapsto Q_m=\tfrac{c}{N}\textstyle\sum_iQ_i$\\
OSF & MiLoRA & A3 & $(U,\Sigma,V)\mapsto\big(U\,\Sigma\,S_{lo}^{-1/2},\ S_{lo}^{1/2}V^\top\big)$\\
P-Tuning & Prompt Tuning & A3 & $h(P,\psi)=f_\psi(P)$\\
PaFi & Diff Pruning & A2 & $x\mapsto(\alpha_0,\ \iota_M x)$\\
Parallel Adapter & AdaptFormer & A3 & $h(W_{down},b_{down},W_{up},b_{up})=\big(W_{down},b_{down},\tfrac{s}{0.1}W_{up},\tfrac{s}{0.1}b_{up}\big)$\\
PEANuT & LoRA & A3 & $h(\Theta_1,\Theta_2)=(\Theta_2^\top,\ \sigma(\Theta_1^\top W_0))$\\
PEM Arithmetic & Task Arithmetic & A3 & $h(\lambda)=(\lambda,\,1-\lambda)$\\
PEQA & EfficientQAT & A3 & $\Delta s\mapsto\Delta s\otimes\mathbf 1_{n/g}^\top$\\
PeriodicLoRA & ReLoRA & A3 & $h=\mathrm{id}\ \text{(stagewise)}$\\
PERK & LoRA & A3 & $h_c(\theta_{adapter},\alpha,\psi)=\phi^{(N)}(\theta_{adapter},\alpha,\psi;c)$\\
PHATGOOSE & LoRA & A3 & $h_x(v)=\big([w_{x,1}(v)B_1\ \cdots\ w_{x,N}(v)B_N],\ [A_1;\dots;A_N]\big),\ \ w_{x,z}=0\ \text{off the top-}k$\\
PHM-Adapter & Adapter (Houlsby) & A3 & $h(\{A_i^D,B_i^D,A_i^U,B_i^U\})=\big(W_{down}=\textstyle\sum_iA_i^D\otimes B_i^D,\;W_{up}=\sum_iA_i^U\otimes B_i^U\big)$\\
Pico & LoRA & A3 & $h(\ast)=\big([\,\gamma SB_1\ \cdots\ \gamma SB_T\,],\ [A_1;\dots;A_T]\big)$\\
Piggyback & Diff Pruning & A3 & $M\mapsto\big(\alpha_0,\ -(\mathbf 1-M)\odot W_0\big)$\\
Piggyback & HiRA & A3 & $M\mapsto(M-\mathbf 1,\ I_n)$\\
PiSSA & LoRA & A3 & $h(B,A)=([B,\,-B_0],\,[A;\,A_0])$\\
PLoP & LoRA & A2 & $h(q_T)=\big(q_T,\ q_0|_{T^{c}}\big)$\\
PoLAR & LoRA & A3 & $h(X,\Theta,Y)=(X\Theta,\ Y^\top)$\\
Polytropon (Poly) & LoRA & A3 & $h_\tau(A,B,Z)=\big(\textstyle\sum_j\alpha_{\tau j}B_j,\ \sum_j\alpha_{\tau j}A_j\big)$\\
Polytropon (Poly) & MHR & A3 & $h(A,B,Z)=(A,\,B,\,Z\otimes\mathbf 1_h)$\\
PPT & Prompt Tuning & A3 & $h=\mathrm{id}_{\mathbb R^{\ell\times d}}$\\
Prefix-Tuning & Initial State Tuning & A3 & $h\big(x^{(l)}_{-\ell..-1}\big)=\Big(\textstyle\sum_{i=1}^{\ell}\bar A_l^{\,i-1}\bar B_l\,x^{(l)}_{-i}\Big)_{l}$\\
Prefix-Tuning & MAM Adapter & A5 & $h(P',\theta)=\big(P',\,\theta,\,W_{\mathrm{up}}=0,\,W^{0}_{\mathrm{down}}\big)$\\
Prefix-Tuning & P-Tuning v2 & A3 & $h(P',\theta)=\big(\mathrm{MLP}_\theta(P'),\,W^{0}_{\mathrm{cls}}\big)$\\
PreFT & LoReFT & A3 & $(W_1,W_2,b)\mapsto\big(R=(W_2W_2^{\top})^{-1/2}W_2,\ W=R+s_r(W_2W_2^{\top})^{1/2}W_1,\ b'=s_r(W_2W_2^{\top})^{1/2}b\big)$\\
PRISM & LoRA & A3 & $(B,A)\mapsto\big([B,\ B_0],\ [A;\ -A_0]\big)$\\
PRISM & PiSSA & A3 & $(B,A)\mapsto(\sqrt s\,B,\ \sqrt s\,A)$\\
PRoLoRA & LoRA & A3 & $h(B_u,B_0,A_u,A_0)=\big(\mathcal B(B_u,B_0),\;\mathcal A(A_u,A_0)\big)$\\
Prompt Tuning & ATTEMPT & A3 & $h(P)=\Big(P_{\mathrm{tgt}}=\tfrac{(t+1)P-\sum_{j\le t}P_j}{t+2},\ \gamma=0,\ \beta=0,\ W^0_{\mathrm{down}},\ W^0_{\mathrm{up}}\Big)$\\
Prompt Tuning & MPT & A3 & $h(P)=\big(P\oslash u_0v_0^\top,\,u_0,\,v_0\big)$\\
Prompt Tuning & P-Tuning & A3 & $h(P)=\big(P'_0,\;\psi_0\ \text{with}\ W_2\mapsto W_2^0+H_0^{+}(P-P_0)\big)$\\
Prompt Tuning & P-Tuning v2 & A3 & $h(p)=\big((K_j(p),V_j(p))_{j\le L},\,W^0_{\mathrm{cls}}\big)$\\
Prompt2Effect & LoRA & A3 & $h_c(\phi)=\big(\hat B_\phi(c,W_0),\ \hat A_\phi(c,W_0)\big)$\\
Proxy-Tuning & EFT & A3 & $h=\mathrm{id}$\\
PSOFT & LoRA-XS & A3 & $(Q,\alpha,\beta)\mapsto\Sigma_r\big(D_\beta C(Q)^\top D_\alpha-I\big)\Sigma_r^{-1}$\\
PSOFT & PiSSA & A3 & $(Q,\alpha,\beta)\mapsto\big(U_r\Sigma_rD_\beta C(Q)^\top D_\alpha\Sigma_r^{-1/2},\ \Sigma_r^{1/2}V_r^\top\big)$\\
Punica & LoRA & A3 & $h_u(\ast)=(B_u,\ A_u)$\\
Punica & S-LoRA & A3 & $h=\mathrm{id}_{\prod_uQ_{\mathrm{LoRA}_{r_u}}}$\\
PVeRA & LoRA-FA & A3 & $h(b,d_\mu,d_\sigma)=\alpha\Lambda_bB\Lambda_{d_\mu}$\\
PVeRA & VeRA & A3 & $h(b,d_\mu,d_\sigma)=(b,\ d_\mu)$\\
QA-LoRA & QLoRA & A3 & $(B,A)\mapsto(\tfrac{s}{s^\prime}B,\ AP)$\\
QLoRA & IR-QLoRA & A2 & $(B,A)\mapsto(\tfrac{s}{\alpha}B,\ A,\ 0,\ 0)$\\
Quantum-PEFT & AdaLoRA & A3 & $h(\theta_U,\theta_V,\lambda)=\big(P=Q_P(\theta_U)I_{d_{out}\times K},\;\Lambda=\mathrm{diag}(\lambda),\;Q=(Q_P(\theta_V)I_{d_{in}\times K})^\top\big)$\\
Quantum-PEFT & LoRA & A3 & $h=\big(U(\theta_U)\,\mathrm{diag}(\lambda),\;V(\theta_V)^\top\big)$\\
QZO & EfficientQAT & A3 & $h=\mathrm{id}:\ S\mapsto S$\\
RandLoRA & LoRA & A3 & $h(\Lambda,\Gamma)=\big([B_1\Lambda_1,\dots,B_n\Lambda_n],\ [A\Gamma_1;\dots;A\Gamma_n]\big)$\\
RaSA & LoRA & A3 & $h=\big([\tilde B_i\ \ B_S]\,D_i,\;[\tilde A_i;\,A_S]\big)_i$\\
RED & SSF & A2 & $(\ell_s,\ell_b)\mapsto\big((\gamma,\beta)_{\rm FFN\ out}=(\ell_s,\ell_b);\ \text{all other SSF wrappers at }(\mathbf 1,0)\big)$\\
ReLoRA & COLA & A3 & $h=\mathrm{id}\ \text{(stagewise)}$\\
ReLoRA & LoRA & A3 & $h((B_k,A_k)_k)=([B_1\cdots B_K],\,[A_1;\dots;A_K])$\\
ReLoRA & ReLoRA & A5 & $(q_1,\dots,q_K)\mapsto(q_1,\dots,q_K,\,q_0^{(K+1)})$\\
RepAdapter & LoRA & A3 & $h(W_u,W_d)=\big((sW_0^{(x)}W_u,\ W_d)\big)_{x\in\{q,k,v\}}$\\
Residual Adapters & LoRA-FA & A3 & $h(\alpha,D_1,D_2)=D_2(I+\alpha D_1)-I\ \ \text{as the trainable }B,\ \text{with }A_0=W_0,\ r=C$\\
Residual Prompt Tuning & Prompt Tuning & A3 & $h(P,W_{\rm down},W_{\rm up},\mathrm{LN})=\mathrm{LN}\big(\mathrm{ReLU}(PW_{\rm down})W_{\rm up}\big)+P$\\
Riemannian LoRA & LoRA & A3 & $h=\mathrm{id}_Q$\\
Riemannion & LoRA & A3 & $X\mapsto\big([\,X(B_0^{+}X)^{+},\ -B_0\,],\ [\,B_0^{+}X;\ A_0\,]\big),\qquad X_0=B_0A_0$\\
RotMoLE & LoRA & A3 & $h_x\big((B_i,A_i,q_i)_i,W_g,W_\theta\big)=\big(\,[\,\mathbb 1_{i\in\mathrm{TopK}(x)}\,g_i(x)\,B_iR_i(x)\,]_{i=1}^{n},\ [A_1;\dots;A_n]\big)$\\
rsLoRA & LoRA & A3 & $h(B,A)=(\sqrt r\,B,\ A)$\\
S-LoRA & LoRA Soups (CAT) & A3 & $h_u(\ast)=e_u$\\
S-LoRA & Punica & A3 & $h=\mathrm{id}_{\prod_uQ_{\mathrm{LoRA}_{r_u}}}$\\
S\textsuperscript{2}FT & Diff Pruning & A2 & $x\mapsto(\alpha_0,\ \iota_M x)$\\
S\textsuperscript{2}FT & LoRA & A2 & $V\mapsto(s^{-1}V,\ U_S^\top)$\\
S\textsuperscript{2}FT & LoRA-FA & A3 & $V\mapsto s^{-1}V\quad(A_0:=U_S^\top)$\\
SAM (sparse) & Diff Pruning & A2 & $x\mapsto(\alpha_0,\ \iota_M x)$\\
ScaLoRA & LoRA & A3 & $h(B,A)=\big([\,B,\ -\widetilde B_t\,],\ [\,A;\ \widetilde A_t\,]\big)$\\
SD-LoRA & LoRA & A3 & $h=\big([\alpha_1\hat B_1\ \cdots\ \alpha_{t-1}\hat B_{t-1}\ \ \alpha_tB_t/\lVert B_tA_t\rVert_F],\ [\hat A_1;\dots;\hat A_{t-1};A_t]\big),\ \ D_k=\hat B_k\hat A_k$\\
SEAL & LoRA & A3 & $h=\mathrm{id}$\\
ShareLoRA & LoRA & A3 & $h\big((B_\ell)_\ell,A\big)=(B_\ell,A)_\ell$\\
SHINE & LoRA & A3 & $h_c(\Delta_{\rm meta},m,\psi)=\mathrm{M2P}_\psi\big(\mathcal M(c)\big)_l$\\
SHiRA & RoSA & A5 & $h(\Delta)=\big(B=0,\;A=A_0,\;S=M\odot\Delta\big)$\\
SIFT & Diff Pruning & A2 & $x\mapsto(\alpha_0,\ \iota_M x)$\\
SingLoRA & LoRA & A3 & $h(A)=\big(\tfrac\alpha r\,u\,A_{[1:d_{out}]},\ A_{[1:d_{in}]}^\top\big)$\\
SLAO (Merge before Forget) & LoRA & A3 & $h=\mathrm{id}_Q$\\
SLTrain & RoSA & A3 & $h=\mathrm{id}$\\
SmartFed (MoRE + EEQA) & LoRA & A3 & $h_x(W_{router})=\big(\widehat B\,\mathrm{diag}(\tilde g_x(W_{router})),\ \widehat A\big)$\\
SMEAR & Adapter (Houlsby) & A3 & $h_v(\theta_1,\dots,\theta_N,W_r)=\textstyle\sum_iR(v)_i\theta_i$\\
SMoA & HiRA & A3 & $(B_k,A_k)_k\mapsto\big(P_{out}^\top\,\mathrm{blkdiag}(B_1,\dots,B_K),\ \mathrm{blkdiag}(A_1,\dots,A_K)\,P_{in}\big)$\\
SMT & Diff Pruning & A2 & $x\mapsto(\alpha_0,\ \iota_M x)$\\
SoRA & LoRA & A3 & $h(B,g,A)=(B\,\mathrm{diag}(g),\ A)$\\
Sparse Memory Finetuning & FISH Mask & A3 & $h=\mathrm{id},\ \ M=\mathrm{rows}(\mathcal T_t)$\\
SparseLoRA & LoRA & A3 & $h=\mathrm{id}$\\
Spectral Adapter\textsuperscript{A} & PiSSA & A3 & $(A_U,A_V)\mapsto\big((U_1+A_U)S_1^{1/2},\ S_1^{1/2}(V_1+A_V)^\top\big)$\\
Spectral Adapter\textsuperscript{R} & LoRA-XS & A3 & $(Q_U,Q_V)\mapsto C(Q_U)\,S_1\,C(Q_V)^\top S_1^{-1}-I$\\
Spectral Adapter\textsuperscript{R} & PiSSA & A3 & $(Q_U,Q_V)\mapsto\big(U_1C(Q_U)S_1C(Q_V)^\top S_1^{-1/2},\ S_1^{1/2}V_1^\top\big)$\\
SpIEL & Diff Pruning & A2 & $(\eta,\phi)\mapsto\big(\alpha_0,\ \mathrm{scatter}(\eta,\phi)\big)$\\
SPoT & Prompt Tuning & A3 & $h=\mathrm{id}_{\mathbb R^{\ell\times d}}$\\
SPT & RoSA & A2 & $\big((B_s,A_s)_{s\in S},\,(x_u)_{u\in U}\big)\mapsto\big((B_s,A_s)_{s\in S},\,(0,A_{0,u})_{u\in U};\ \iota(x)\big)$\\
SRR (Structured Residual Reconstruction) & LoRA & A3 & $(B,A)\mapsto\big([B,\ -B_0],\ [A;\ A_0]\big)$\\
SSF & GLoRA & A3 & $(\gamma,\beta)\mapsto\big(B_A=\gamma-\mathbf 1,\ A_A=\mathbf 1^\top;\ D=\gamma-\mathbf 1;\ E=\beta\big)\ \text{(other paths at zero)}$\\
Stable-LoRA & LoRA & A3 & $h=\mathrm{id}_Q$\\
StructLoRA & LoRA & A3 & $(A,B,m)\mapsto\big(\tfrac{\alpha}{r}B\,\mathrm{diag}(m),\ A\big)$\\
StructLoRA & SoRA & A3 & $(A,B,m)\mapsto\big(B,\ \tfrac{\alpha}{r}m,\ A\big)$\\
Super-Tuning (Super / Supra) & PaFi & A5 & $h(U)=\big(U,\,\gamma_0,\,\beta_0\big)$\\
Super-Tuning (Super / Supra) & SHiRA & A3 & $h=\mathrm{id}$\\
Surgical FT & Diff Pruning & A2 & $x\mapsto(\alpha_0,\ \iota_M x)$\\
Surgical FT & HiFT & A5 & $q\mapsto(\text{stage } j: q;\ \text{other stages}: 0)$\\
SVDiff & SVF (Transformer\textsuperscript{2}) & A3 & $\delta\mapsto\mathrm{ReLU}(\sigma+\delta)\oslash\sigma$\\
SVDiff & SVFT & A3 & $\delta\mapsto\mathrm{ReLU}(\sigma+\delta)-\sigma$\\
SVF (Transformer\textsuperscript{2}) & SVFT & A3 & $z\mapsto\sigma\odot(z-\mathbf 1)$\\
SVFT & LoRA-FA & A3 & $h(M)=U_{:,R}\,M_{R,C},\ \ A_0=V_{:,C}^\top$\\
Task Arithmetic & AdaMerging & A3 & $h(\lambda)=(\lambda_i^{l}=\lambda_i)_{i,l}$\\
Task Arithmetic & LoRA Soups (CAT) & A3 & $h(\lambda)=(\alpha_i^{l}=\lambda_i)_{i,l}$\\
Task Arithmetic & ZipLoRA & A3 & $h(\lambda_c,\lambda_s)=(\lambda_c\mathbf 1,\ \lambda_s\mathbf 1)$\\
TeRA & VeRA & A3 & $h(d^{(1)},d^{(2)})=\big(b=d^{(1)},\ d=d^{(2)}\big)$\\
Text-to-LoRA (T2L) & LoRA & A3 & $h_z(\theta)=h_\theta\big(\mathrm{concat}[f(z),E[m],E[l]]\big)$\\
Textual Inversion & Custom Diffusion & A2 & $h(v)=(\Delta W^k=0,\ \Delta W^v=0,\ V^*=v)$\\
Textual Inversion & Trainable Tokens & A3 & $h=\mathrm{id}:\ v\mapsto V_{\{S_*\}}$\\
Tied-LoRA & LoRA & A3 & $h=\big(\Lambda_{v_\ell}B\Lambda_{u_\ell},\;A\big)_\ell$\\
Tied-LoRA & ShareLoRA & A3 & $h(\{v_\ell,u_\ell\}_\ell,B,A)=\big(\{\Lambda_{v_\ell}B\Lambda_{u_\ell}\}_\ell,\;A\big)$\\
Tina & LoRA & A3 & $h=\mathrm{id}$\\
TinyLoRA & LoRA-XS & A3 & $h(v)=\Sigma_r^{1/2}\Big(\sum_iv_iP_i\Big)\Sigma_r^{-1/2}$\\
TinyLoRA & PiSSA & A3 & $h(v)=\Big(U_r\Sigma_r^{1/2}\big(I+\textstyle\sum_iv_iP_i\big),\ \Sigma_r^{1/2}V_r^\top\Big)$\\
Trainable Tokens & LoRA & A2 & $h(X)=(S_T,\,X),\quad X=V_T-E_{0,T}$\\
Trainable Tokens & SHiRA & A3 & $h=\mathrm{id}$\\
Trans-LoRA & LoRA & A3 & $h=\mathrm{id}$\\
TT-LoRA & LoRA & A3 & $h(\mathcal C_1,\dots,\mathcal C_D)=\big(\mathcal C_1\times^1\cdots\times^1\mathcal C_c,\;\mathcal C_{c+1}\times^1\cdots\times^1\mathcal C_D\big)\ \text{reshaped to }d_{out}\times r_c,\ r_c\times d_{in}$\\
Uni-LoRA & LoRA & A3 & $h(\theta_d)=P\theta_d=\big(B^\ell(P\theta_d),A^\ell(P\theta_d)\big)_\ell$\\
VAPT & VPT & A3 & $h_x(q)=\big((P^{(l)}(\tilde X^{(l)}(x;q)))_{l\le L},\ \mathrm{Head}\big)$\\
VB-LoRA & LoRA & A3 & $h(\alpha,\sigma)=\big(B(\alpha,\sigma),\ A(\alpha,\sigma)\big),\ \ u=\textstyle\sum_s\mathrm{softmax}(\mathrm{TopK}(\sigma,k))_s\,\alpha_s$\\
VeLoRA & LoRA & A3 & $h=\mathrm{id}_Q$\\
VeRA & LoRA & A3 & $h(b,d)=(\Lambda_bB\Lambda_d,\ A)$\\
VeRA & LoRA-FA & A3 & $h(b,d)=\Lambda_bB\Lambda_d$\\
VPT & P-Tuning v2 & A3 & $h(P,\mathrm{Head})=\Big(\big(\mathrm{LN}(P_{i-1})W_k^{(i)}+\mathbf 1b_k^{(i)\top},\,\mathrm{LN}(P_{i-1})W_v^{(i)}+\mathbf 1b_v^{(i)\top}\big)_{i\le L},\,\mathrm{Head}\Big)$\\
WaveFT & LoRA-FA & A3 & $h(c)=\lambda\,B_W\,C_\Omega(c),\ \ A_W=[\psi_j^\top]_{j\in\mathrm{cols}(\Omega)}$\\
WiSE-FT & Concept Sliders & A3 & $h(\alpha)=(\alpha B_1,\ A_1)$\\
WiSE-FT & Task Arithmetic & A3 & $h(\alpha)=(\lambda=\alpha)$\\
X-LoRA & LoRA Soups (CAT) & A3 & $h_x(\phi)=\big(\alpha_i^{l}=\lambda_i^{l}(x;\phi)\,\alpha_i/r_i\big)$\\
XPrompt & Prompt Tuning & A2 & $h(P_S)=\big(P_S,\,0_{S^c}\big)$\\
ZipLoRA & LoRA & A3 & $h(m_c,m_s)=\big([B_c\ B_s],\ [A_c\,\mathrm{diag}(m_c);\ A_s\,\mathrm{diag}(m_s)]\big)$\\
ZO-Muon & LoRA & A2 & $X\mapsto(P_k,\ X)$\\
\end{longtable}\endgroup

